\documentclass[11pt]{article}
\usepackage{tcee}
\addbibresource{3A17.bib}
\tcetema[Análisis de mercados (II): el monopolio]{3.A.17}{Análisis de mercados (II). Teoría del monopolio. Discriminación de precios. Monopolio natural. Producción conjunta. Análisis de eficiencia y bienestar. Monopsonio. Monopolio bilateral.}
\tcefecha{17/10/2026}
\begin{document}
\tceetitulo

\section*{Introducción}
\minutos{2,5}

\begin{itemize}
\item \textbf{\emph{Enganche}}:
  \begin{itemize}
  \item «Quien parte y reparte se queda con la mejor parte», dice el refrán. Una empresa sin competidores parte y reparte: decide cuánto vende y a qué precio. El 5 de agosto de 2024, un juez federal de Estados Unidos concluyó que «Google es un monopolista y ha actuado como tal para mantener su monopolio», con una cuota del 89,2~\% en las búsquedas generales en internet y del 94,9~\% en los teléfonos móviles \parencite{Mehta2024Google}.
  \item Detrás está la pregunta de esta exposición: ¿se queda siempre con la mejor parte quien es el único que vende? ¿Qué le limita? ¿Y qué le queda al resto de la sociedad?
    \begin{itemize}
    \item La respuesta de la microeconomía parte del \concepto{individualismo metodológico}:\footnote{Es un método muy extendido en las ciencias sociales: sostiene que los fenómenos sociales, su estructura y sus cambios, se explican en principio a partir de los individuos, de sus metas, sus creencias y sus acciones, y ve la evolución de la sociedad como el agregado de las decisiones de los particulares. Es, por tanto, un reduccionismo: explica las entidades grandes a partir de las pequeñas.} explica los fenómenos agregados a partir de las decisiones de individuos que optimizan \vertema{3.A.8}.
    \end{itemize}
  \item Para comprender y predecir el funcionamiento de los mercados, la \concepto{microeconomía} estudia 2 agentes, \emph{consumidores} y \emph{productores}:\footnote{La microeconomía contemporánea considera esta separación estricta entre consumidores y productores una simplificación excesiva \parencite[cap.~26]{EkelundHebert2013History}: lo muestran las «tecnologías del consumo», como el enfoque de características de \autor{Lancaster}, la economía doméstica de \autor{Becker}, la producción doméstica de \autor{Gronau} o la economía de la información de \autor{Stigler}, en la que conocer los precios y las calidades exige una búsqueda costosa que se combina con el bien físico. La microeconomía estudia también otros agentes, como el sector público o las instituciones financieras.} los consumidores eligen qué consumir sujetos a su restricción presupuestaria, y los productores transforman factores (\emph{inputs}) en productos (\emph{outputs}). Las decisiones de los productores están restringidas por aspectos:
    \begin{itemize}
    \item \emph{técnicos}: cómo se combinan los factores de forma eficiente, dada una tecnología, que estudia la teoría de la producción \vertema{3.A.11};
    \item \emph{de costes}: cuál de las combinaciones técnicamente eficientes es la más barata, que estudia la teoría de los costes \vertema{3.A.12};
    \item \emph{organizativos y de mercado}: el tamaño y los límites de la empresa \vertema{3.A.15} y el mercado en el que vende y compra \vertema{3.A.16} \vertema{3.A.18} \vertema{3.A.19}.
    \end{itemize}
  \item El resultado de la interacción de los agentes, el \concepto{equilibrio de mercado}, depende de la \concepto{estructura de mercado}. Se suelen distinguir 2 estructuras polares, la competencia perfecta y el monopolio, y 2 intermedias, el oligopolio y la competencia monopolística.\footnote{En realidad no puede afirmarse que existan solo 4 estructuras de mercado: la clasificación facilita la exposición, y cada una tiene su tema \vertema{3.A.16} \vertema{3.A.18} \vertema{3.A.19}. En la práctica pueden mencionarse otras, como las \emph{plataformas}, que incorporan las economías de red: la utilidad de cada usuario aumenta con el número de usuarios.} Estudiamos el monopolio en \emph{equilibrio parcial}.\footnote{Es el enfoque de \textcite{Marshall1890Principles}: se estudia el mercado de un único bien, que es una parte pequeña de la economía, de modo que lo que ocurre en él no altera los precios ni las rentas del resto \vertema{3.A.16}.}
  \item El \concepto{monopolio} (del griego \emph{mónos}, uno, y \emph{pōleîn}, vender) es la estructura de mercado con un único vendedor de un bien sin sustitutivos cercanos, protegido por barreras de entrada. Su rasgo esencial es el \concepto{poder de mercado}: la capacidad de fijar un precio por encima del coste marginal sin perder todas las ventas. Como la empresa deja de ser precio-aceptante, falla un supuesto del primer teorema fundamental de la economía del bienestar (1TFEB), y el equilibrio deja de ser eficiente \vertema{3.A.22}.
  \item Su simétrico en el lado de la compra es el \concepto{monopsonio}, un único comprador, y, cuando un monopolista se enfrenta a un monopsonista, hablamos de \concepto{monopolio bilateral}.
  \end{itemize}

\item \textbf{\emph{Relevancia}}:
  \begin{itemize}
  \item \emph{Teórica}: el monopolio es el caso más sencillo de poder de mercado, y sus herramientas (el ingreso marginal, el índice de \autor{Lerner}, la autoselección) sirven para todo el análisis de los mercados imperfectos y de la regulación.
  \item \emph{Empírica}: el poder de mercado ha vuelto al centro del debate. Con datos de empresas cotizadas de Estados Unidos, el precio medio pasó de estar un 21~\% por encima del coste marginal en 1980 a un 61~\% en 2016 \parencite{DeLoeckerEeckhoutUnger2020Rise}, y el poder de las empresas como compradoras de trabajo es apreciable \parencite{YehMacalusoHershbein2022Monopsony}.
  \item \emph{De política económica}: es la base teórica de la defensa de la competencia \vertema{3.B.39}, de la regulación de los monopolios naturales \vertema{3.A.20} y del debate sobre el salario mínimo cuando el empleador tiene poder de mercado \vertema{3.A.25}.
  \end{itemize}

\item \textbf{\emph{Contextualización}}. Desde un punto de vista histórico:
  \begin{itemize}
  \item En el siglo XVIII, Adam \textcite{Smith1776Wealth} ya observa que el monopolio concedido a un individuo o a una compañía comercial mantiene el mercado desabastecido para vender muy por encima del precio natural: «el precio de monopolio es, en toda ocasión, el más alto que se puede obtener» (libro~I, cap.~VII).
  \item En el siglo XIX, David \textcite{Ricardo1817Principles} define el precio de monopolio como el más alto que los consumidores están dispuestos a pagar, sin relación con el coste de producción; John Stuart \textcite{Mill1848Principles} estudia el valor de los productos conjuntos (el coque y el gas), y Jules \textcite{Dupuit1849Peages} analiza por primera vez la discriminación de precios con los peajes de las vías de comunicación \parencite{Ekelund1970Price}, tras medir en \textcite{Dupuit1844Mesure} la utilidad que se pierde cuando el precio excede el coste, el germen de la pérdida irrecuperable.
  \item La primera solución matemática del monopolio es la de \textcite{Cournot1838Recherches}, con el dueño de un manantial de agua mineral y la condición de que el ingreso marginal iguale el coste marginal. Le siguen \autor{Bertrand} (1883), \autor{Edgeworth} (1897), \textcite{Hotelling1929Stability} y \textcite{Robinson1933Economics}. Estos modelos toman el poder de mercado como dado y no prestan mucha atención a sus causas ni a la entrada de competidores.\footnote{Los modelos formales de los mercados con pocas empresas suelen suponer, siguiendo el ejemplo de \autor{Cournot}, que la entrada es imposible; lo advierte \textcite{Modigliani1958New}, en el origen de la teoría de la disuasión de la entrada \vertema{3.A.15}.}
  \item \textcite{Marshall1890Principles} formaliza el excedente del consumidor y el análisis del monopolio en equilibrio parcial, y su sucesor en Cambridge, Arthur C. \textcite{Pigou1920Economics}, clasifica la discriminación de precios en 3 grados. Joan \textcite{Robinson1933Economics} populariza el ingreso marginal, desarrolla la discriminación de tercer grado e introduce el término \emph{monopsonio}; Arthur \textcite{Bowley1928Bilateral} estudia el monopolio bilateral, y Abba \textcite{Lerner1934Concept} propone el índice de poder de mercado que lleva su nombre.
  \item Tras la Segunda Guerra Mundial, \textcite{Harberger1954Monopoly} cuantifica la pérdida de bienestar del monopolio; \textcite{Coase1946Marginal} y \textcite{Boiteux1956Gestion} abren la teoría de la tarificación de los monopolios públicos; \textcite{Baumol1977Proper} define el monopolio natural por la subaditividad de los costes, y la nueva organización industrial incorpora la información asimétrica a la discriminación de precios \parencites{MussaRosen1978Monopoly}{MaskinRiley1984Monopoly} y a la regulación \parencite{LaffontTirole1993Theory}.
  \end{itemize}

\item \textbf{\emph{Problemática (preguntas clave)}}:
  \begin{itemize}
  \item ¿Cómo fija el precio y la cantidad una empresa sin competidores, qué limita su poder y cuánto le cuesta eso a la sociedad?
  \item ¿Qué cambia cuando puede cobrar precios distintos, cuando la tecnología hace preferible una sola empresa o cuando vende varios productos?
  \item ¿Y si el poder de mercado está del lado de quien compra, o en los 2 lados a la vez?
  \end{itemize}

\item \textbf{\emph{Estructura}}: la exposición se divide en 3 bloques.
  \begin{itemize}
  \item \concepto{I. Teoría del monopolio}: el escenario base (una sola empresa, muchos compradores precio-aceptantes y barreras a la entrada) (\ref{sec:caracterizacion-monopolio-escenario}), su equilibrio y la regla de \autor{Lerner} (\ref{sec:equilibrio-monopolio}) y el análisis de eficiencia y bienestar (\ref{sec:analisis-eficiencia-bienestar}), que adelantamos porque es la vara de medir de todo lo que sigue.
  \item \concepto{II. Desarrollos de la teoría del monopolio}, levantando de uno en uno los supuestos del escenario base: el precio único (discriminación de precios, \ref{sec:discriminacion-precios}), unos costes sin ventajas de tamaño (monopolio natural y su regulación, \ref{sec:monopolio-natural}) y el producto y la planta únicos (producción conjunta, \ref{sec:produccion-conjunta}).
  \item \concepto{III. Otras estructuras de mercado relacionadas con el monopolio}: el poder de mercado cambia de lado (monopsonio, \ref{sec:monopsonio}) y, por último, está en los 2 lados a la vez (monopolio bilateral, \ref{sec:monopolio-bilateral}).
  \end{itemize}
  Un mismo hilo recorre los 3 bloques: la \concepto{regla de la elasticidad inversa}, que reaparece por grupos de consumidores, amortiguada por una restricción presupuestaria, con elasticidades cruzadas y, en espejo, en el mercado de factores.
\end{itemize}

\begin{esquema}
\textbf{3.A.17. Monopolio, monopsonio y monopolio bilateral}\\[2pt]
Introducción\\
I. Teoría del monopolio\\
\hspace*{1.5em}1. Escenario base: 1 oferente, barreras\\
\hspace*{1.5em}2. Equilibrio: $\mathrm{IMg}=\mathrm{CMg}$ $\Rightarrow$ $\mathcal{L}=1/|\varepsilon_{q,p}|$\\
\hspace*{1.5em}3. Eficiencia y bienestar: \autor{Harberger}, \autor{Tullock}-\autor{Posner}, ineficiencia X\\
II. Desarrollos de la teoría del monopolio\\
\hspace*{1.5em}1. Discriminación de precios: 1.\textsuperscript{er}, 2.º (autoselección), 3.\textsuperscript{er} grado, intertemporal\\
\hspace*{1.5em}2. Monopolio natural: subaditividad, sostenibilidad, regulación\\
\hspace*{1.5em}3. Producción conjunta: multiproducto, proporciones fijas, multiplanta\\
III. Otras estructuras relacionadas con el monopolio\\
\hspace*{1.5em}1. Monopsonio: $\mathrm{VPMg}_L=\mathrm{GMg}_L$\\
\hspace*{1.5em}2. Monopolio bilateral: indeterminación $\rightarrow$ negociación de \autor{Nash}\\
Conclusión
\end{esquema}

\begin{notaopositor}
La exposición va del monopolio básico y su coste social (\ref{sec:teoria-monopolio-equilibrio}) a los desarrollos que el título pide (\ref{sec:desarrollos-teoria-monopolio}: discriminación de precios, monopolio natural y producción conjunta) y a las estructuras en que el poder de mercado está en el lado de la demanda o en los 2 lados (\ref{sec:otras-estructuras-mercado}). Reparto del cante: introducción, \tiempo{2,5}; I, \tiempo{8} (I.1, \tiempo{1}; I.2, \tiempo{4}; I.3, \tiempo{3}); II, \tiempo{11,5} (II.1, \tiempo{5,5}; II.2, \tiempo{4,5}; II.3, \tiempo{1,5}); III, \tiempo{7} (III.1, \tiempo{3,5}; III.2, \tiempo{3,5}), y conclusión, \tiempo{1}. El análisis de eficiencia y bienestar va en \ref{sec:analisis-eficiencia-bienestar}, antes que los desarrollos, aunque en el título va detrás de la producción conjunta: es la comparación del equilibrio de \ref{sec:equilibrio-monopolio} con el competitivo y lo usan todos los desarrollos, que terminan con su propio juicio de bienestar. Los gráficos imprescindibles son el equilibrio del monopolio con sus capas de bienestar (triángulo de \autor{Harberger}, búsqueda de rentas e ineficiencia X), la autoselección con 2 tipos, la subaditividad y la sostenibilidad, la regulación del monopolio natural y el monopsonio, sobre el que se construye el monopolio bilateral. La función clave del índice de \autor{Lerner} se escribe en \ref{sec:equilibrio-monopolio} y se queda en una esquina de la pizarra todo el cante. La regulación con información asimétrica (tasa de rentabilidad, \emph{price cap}, menú de contratos, competencia por comparación) y la teoría positiva de la regulación están en el anexo~\ref{anexo:regulacion}, para las preguntas del tribunal \vertema{3.A.20}. Conviene escribir este esquema en la pizarra al empezar.
\end{notaopositor}


% ==========================================================================
\section{Teoría del monopolio: equilibrio,\newline eficiencia y bienestar}\label{sec:teoria-monopolio-equilibrio}
\minutos{8}

\begin{notaopositor}
En el cante: el escenario base en una lista corta, con la tabla de supuestos como apoyo (\tiempo{1}); el problema del monopolista, la condición $\mathrm{IMg}=\mathrm{CMg}$, el índice de \autor{Lerner} como función clave, el tramo elástico, la inexistencia de la curva de oferta y la condición de viabilidad, con el gráfico del equilibrio (\tiempo{4}); y el bienestar sobre el mismo gráfico: transferencia, triángulo de \autor{Harberger} y su tamaño, búsqueda de rentas e ineficiencia X, con la evidencia en una frase (\tiempo{3}). La condición de segundo orden, el corto plazo y el largo plazo en detalle, los impuestos y la eficiencia dinámica son de profundización.
\end{notaopositor}

\begin{itemize}
\item Empezamos por el modelo básico. Su \emph{conclusión anticipada} es sencilla: el monopolista produce menos y cobra más que la competencia, el margen que puede cargar sobre el coste marginal lo limita la elasticidad de la demanda, y la pérdida de bienestar que causa va más allá de las unidades que deja de producir.
\end{itemize}

% --------------------------------------------------------------------------
\subsection{Caracterización del monopolio: el escenario base}\label{sec:caracterizacion-monopolio-escenario}

\begin{itemize}
\item Empezamos por el caso más sencillo, que servirá de \concepto{escenario base}. Parte de la misma lista de supuestos que todas las estructuras de mercado (tabla~\ref{tab:supuestos}) y \emph{levanta el supuesto~4}: hay un solo vendedor, que fija el precio; y, en lugar de la libre entrada del supuesto~5, hay barreras de entrada.
\end{itemize}

\begin{table}[H]
\centering
\caption{Supuestos de las estructuras de mercado: el monopolio levanta el 4 (y el 5)}\label{tab:supuestos}
\begin{tabular}{>{\centering\arraybackslash}m{0.7cm}>{\raggedright\arraybackslash}m{4.3cm}>{\raggedright\arraybackslash}m{10.0cm}}
\toprule
\rowcolor{tceeazulpalido}\textbf{N.º} & \textbf{Supuesto} & \textbf{Qué significa}\\
\midrule
\textcolor{tceemorado}{\textbf{1}} & \textbf{Agentes racionales} & los consumidores maximizan su utilidad y las empresas, su beneficio\\
\addlinespace[0.3em]
\textcolor{tceemorado}{\textbf{2}} & \textbf{Información perfecta y simétrica} & los consumidores conocen las características del producto y el precio, y las empresas, su tecnología y el precio; nadie sabe más que los demás\\
\addlinespace[0.3em]
\textcolor{tceemorado}{\textbf{3}} & \textbf{Producto homogéneo (sin sustitutivos cercanos)} & los consumidores ven como sustitutivos perfectos las unidades de todas las empresas, de modo que solo les importa el precio\\
\addlinespace[0.3em]
\rowcolor{tceesalmon}\textcolor{tceemorado}{\textbf{4}} & \textbf{Mercado atomizado: muchos compradores y vendedores (precio-aceptantes)} & cada uno es pequeño respecto al mercado y sus decisiones individuales no mueven el precio: todos son \concepto{precio-aceptantes}, y el precio lo determinan las decisiones de todos juntos, no las de ninguno en particular\newline \textcolor{tceerojo}{\textbf{Se levanta}}: hay \textbf{un solo vendedor}, que fija el precio (\emph{precio-decisor}); los compradores siguen siendo muchos y precio-aceptantes\\
\addlinespace[0.3em]
\rowcolor{tceesalmon}\textcolor{tceemorado}{\textbf{5}} & \textbf{Libre entrada y salida (a largo plazo)} & a largo plazo, si hay beneficios extraordinarios, entran empresas hasta agotarlos, y si hay pérdidas, salen\newline \textcolor{tceerojo}{\textbf{Se levanta}}: hay \textbf{barreras de entrada}, que protegen los beneficios del monopolista también a largo plazo\\
\bottomrule
\end{tabular}
\fuente{elaboración propia.}
\end{table}

\begin{itemize}
\item Supuesto a supuesto, el escenario base del monopolio queda así:
  \begin{enumerate}
  \item \emph{Se mantiene: agentes racionales}. Los consumidores maximizan su utilidad y el monopolista, su beneficio.
  \item \emph{Se mantiene: información perfecta y simétrica}. El monopolista conoce la demanda y su tecnología, y los consumidores conocen las características del producto y el precio.
  \item \emph{Se mantiene: un único bien, sin sustitutivos cercanos}. La elasticidad cruzada de su demanda con los precios de los bienes de otras industrias es pequeña, de modo que el monopolista no teme perder clientes hacia ellas.\footnote{Que haya \emph{algún} efecto cruzado no elimina el poder de mercado: lo reduce si los sustitutivos son cercanos, porque la demanda del monopolista se vuelve más elástica. Si hay muchos vendedores de variedades cercanas, es la competencia monopolística \vertema{3.A.18}.}
  \item \emph{Se levanta: un único oferente y muchos demandantes}. Al ser muchos, los compradores se comportan como precio-aceptantes y no pueden revender el bien entre sí; el único oferente, en cambio, tiene poder de mercado.
  \item \emph{Se levanta: barreras de entrada} que impiden a otras empresas competir \vertema{3.A.15}:
    \begin{itemize}
    \item \emph{estratégicas}: conductas del monopolista que disuaden la entrada de competidores potenciales;
    \item \emph{legales o administrativas}: patentes, licencias, concesiones o aranceles que excluyen a los competidores extranjeros;
    \item \emph{tecnológicas}: una tecnología que hace que los costes sean menores con un único productor (\emph{subaditividad} de costes); es el \emph{monopolio natural} (\ref{sec:monopolio-natural}).
    \end{itemize}
  \end{enumerate}
\item A esa lista el escenario base añade 2 supuestos propios del monopolio, que los desarrollos irán levantando (tabla~\ref{tab:escenario-base}): un \emph{precio uniforme} para todas las unidades y todos los compradores, y una \emph{sola planta}, con los factores comprados a precios dados, en mercados competitivos.
\item La consecuencia de levantar el supuesto~4 es el \concepto{poder de mercado}: el monopolista se enfrenta a toda la demanda del mercado, que tiene pendiente negativa, y por eso es \emph{precio-decisor}: puede fijar un precio por encima del coste marginal sin perder todas sus ventas. Pero elige un punto de esa demanda: el precio o la cantidad, no los 2 por separado.\footnote{El poder de mercado no es, por tanto, total: si fija el precio, la demanda le dice cuánto venderá; si fija la cantidad, la demanda le dice a qué precio.} Y, por las barreras (supuesto~5), puede mantener beneficios extraordinarios a largo plazo.
\end{itemize}

\begin{table}[H]
\centering
\caption{Los desarrollos del monopolio: qué supuesto del escenario base levanta cada uno}\label{tab:escenario-base}
\begin{tabular}{>{\centering\arraybackslash}m{0.42\textwidth}>{\centering\arraybackslash}m{0.42\textwidth}}
\toprule
\rowcolor{tceeazulpalido}\textbf{Supuesto del escenario base} & \textbf{Desarrollo que lo levanta}\\
\midrule
Un único precio para todas las unidades\newline y todos los consumidores & Discriminación de precios (\ref{sec:discriminacion-precios})\\
Información completa sobre\newline la valoración de cada consumidor & Discriminación de segundo grado:\newline autoselección (\ref{sec:discriminacion-segundo-grado})\\
Costes sin ventajas de tamaño\newline frente a la demanda & Monopolio natural (\ref{sec:monopolio-natural})\\
Un único producto y una única planta & Producción conjunta (\ref{sec:produccion-conjunta})\\
Precio-aceptante en el mercado\newline de factores & Monopsonio (\ref{sec:monopsonio})\\
Demandantes sin poder de mercado & Monopolio bilateral (\ref{sec:monopolio-bilateral})\\
\bottomrule
\end{tabular}
\fuente{elaboración propia.}
\end{table}

% --------------------------------------------------------------------------
\subsection{Equilibrio del monopolio}\label{sec:equilibrio-monopolio}

\begin{itemize}
\item Una vez caracterizado el escenario base, pasamos a obtener el equilibrio del monopolista: cuánto produce y a qué precio vende.
\end{itemize}

\paragraph{a) Idea}

\begin{itemize}
\item En competencia perfecta la empresa toma el precio como dado. El monopolista no: sabe que, para vender una unidad más, tiene que bajar el precio de \emph{todas} las que vende. Esa es toda la diferencia, y de ella sale todo lo demás.
\item Su estudio permite, desde un punto de vista \emph{positivo}, determinar el precio y la cantidad que maximizan el beneficio del monopolista y, desde un punto de vista \emph{normativo}, valorar si ese resultado es deseable para la sociedad (\ref{sec:analisis-eficiencia-bienestar}).
\end{itemize}

\paragraph{b) Modelo}

\subparagraph{i) Supuestos}

\begin{itemize}
\item Los del escenario base (\ref{sec:caracterizacion-monopolio-escenario}), con 2 supuestos técnicos:
  \begin{itemize}
  \item la \concepto{demanda de mercado} $x(p)$ es derivable y decreciente. El monopolista vende $q=x(p)$, así que puede trabajar con la \concepto{demanda inversa} $p(q)$, con $p'(q)<0$: el precio al que puede vender cada cantidad. Procede de consumidores que maximizan una utilidad cuasilineal sujeta a su restricción presupuestaria \vertema{3.A.8} \vertema{3.A.16}, y mide la cantidad que desean comprar a cada precio, \emph{ceteris paribus};\footnote{En el siglo XVI, Juan de \autor{Medina} y Luis de \autor{Molina}, de la Escuela de Salamanca, ya usan la cláusula \emph{ceteris paribus} al razonar sobre cuestiones económicas; \autor{Molina} lo hace al tratar el valor del dinero en el tomo~II de \emph{De iustitia et iure} (1597) \parencite{SEP2024Ceteris}. En la misma escuela, Martín de \autor{Azpilcueta} (1556) relaciona la llegada del oro y la plata de América con la subida de los precios 12 años antes que Jean \autor{Bodin} (1568). La cláusula la populariza después \autor{Marshall} con su enfoque de equilibrio parcial.}
  \item la \concepto{función de costes} $C(q)$ es derivable y creciente, $\mathrm{CMg}(q)=C'(q)>0$, y ya recoge la combinación de factores más barata para cada cantidad, con los precios de los factores dados: es $C(\vect{w},q)$ con $\vect{w}$ fijo \vertema{3.A.12}.
  \end{itemize}
\end{itemize}

\subparagraph{ii) Desarrollo}

\begin{itemize}
\item \emph{Problema del monopolista}. Elige la cantidad (o, lo que es equivalente, el precio)\footnote{Es indiferente fijar el precio o la cantidad, porque están relacionados de forma unívoca a través de la función de demanda.} que maximiza su beneficio. A diferencia de la competencia perfecta, el precio es una variable endógena:
  \begin{equation}
  \optim{\max}{q}{B(q)=p(q)\,q-C(q)}{q\geq 0\\ B(q)\geq 0} \label{eq:monopolio}
  \end{equation}
  La segunda restricción es la \concepto{condición de participación}: el monopolista solo produce si no tiene pérdidas; si las tuviera, no entraría en el mercado (o cerraría).
\item \emph{Condición de primer orden}: produce hasta que el ingreso que le da la última unidad iguala lo que le cuesta:\footnote{La condición $\mathrm{IMg}=\mathrm{CMg}$ la cumple cualquier empresa que maximiza el beneficio, sea cual sea la estructura de mercado (en competencia perfecta, $\mathrm{IMg}=p$); \textcite{Cournot1838Recherches} la formula por primera vez para el monopolio.}
  \begin{equation}
  B'(q^M)=0\ \Longrightarrow\ \underbrace{p(q^M)+p'(q^M)\,q^M}_{\mathrm{IMg}(q^M)}=\mathrm{CMg}(q^M). \label{eq:img-cmg}
  \end{equation}
  \begin{itemize}
  \item El \concepto{ingreso marginal} (IMg) tiene 2 partes: el precio que cobra por la unidad adicional, $p$, y lo que pierde en las que ya vendía al tener que bajar el precio, $p'(q)\,q<0$. Si el IMg fuera menor que el CMg, le convendría producir menos, y al revés.
  \item Escrito con la elasticidad de la demanda, $\varepsilon_{q,p}=\frac{dq}{dp}\frac{p}{q}<0$:
    \begin{equation}
    \mathrm{IMg}=p\left(1+\frac{dp}{dq}\frac{q}{p}\right)=p\left(1-\frac{1}{|\varepsilon_{q,p}|}\right)<p . \label{eq:img-elasticidad}
    \end{equation}
    En competencia perfecta, la demanda a la que se enfrenta cada empresa es perfectamente elástica, $|\varepsilon_{q,p}|\to\infty$, y $\mathrm{IMg}=p$. En el monopolio la elasticidad es finita y el IMg queda por debajo del precio:\footnote{Si la demanda del mercado fuera perfectamente elástica, coincidiría con el ingreso marginal y se alcanzaría la solución competitiva, $\mathrm{IMg}=p=\mathrm{CMg}$.} lo que la última unidad aporta al monopolista es menos de lo que vale para los consumidores.
  \end{itemize}
\item \emph{Tramo elástico}: como $\mathrm{CMg}>0$, en el óptimo el IMg también ha de ser positivo, y por \eqref{eq:img-elasticidad} eso exige $|\varepsilon_{q,p}|>1$ (estrictamente): el monopolista nunca produce en el tramo inelástico de la demanda.\footnote{En el tramo inelástico el IMg es negativo: reducir la producción aumenta el ingreso y reduce el coste, así que el beneficio aumenta y no puede estar en un máximo. Por la misma razón, tampoco produce donde $|\varepsilon_{q,p}|=1$ ($\mathrm{IMg}=0$, donde se maximiza el ingreso), porque reducir la producción no cambia el ingreso y reduce el coste. Solo produciría ahí con un coste marginal nulo.\graficonota{tramo-elastico}}
\item \emph{Índice de \autor{Lerner}}. Reordenando la condición de primer orden se obtiene la \concepto{regla de la elasticidad inversa}: el monopolista carga sobre el coste marginal un margen (\emph{mark-up}) que es la inversa de la elasticidad de la demanda en el equilibrio,\footnote{Derivación: $\dfrac{p-\mathrm{CMg}}{p}=\dfrac{p-\mathrm{IMg}}{p}=\dfrac{p-\big(p'(q)\,q+p\big)}{p}=-p'(q)\,\dfrac{q}{p}=\dfrac{1}{|\varepsilon_{q,p}|}$.}
  \funcionclave[eq:lerner]{\mathrm{IMg}(q^M)=\mathrm{CMg}(q^M)\ \Longleftrightarrow\ \mathcal{L}\equiv\frac{p^M-\mathrm{CMg}(q^M)}{p^M}=\frac{1}{|\varepsilon_{q,p}|}}
  El \concepto{índice de \autor{Lerner}} mide el poder de mercado como el margen relativo del precio sobre el coste marginal \parencite{Lerner1934Concept}, y lo usan las autoridades de competencia. Implica que:
  \begin{enumerate}[label=\roman*)]
  \item el poder de mercado del monopolista está limitado por la demanda: si es muy elástica, se reduce drásticamente; con demanda perfectamente elástica, $\mathcal{L}=0$ y volvemos a la competencia perfecta;
  \item cuanto más inelástica es la demanda, es decir, cuanto menos reaccionan los consumidores a una subida del precio, mayor es el poder de mercado y mayor el precio.
  \end{enumerate}
\item \emph{Condición de segundo orden}: $B''(q^M)=\mathrm{IMg}'(q^M)-\mathrm{CMg}'(q^M)<0$, es decir, en el óptimo la pendiente del IMg ha de ser menor que la del CMg.\footnote{Las pendientes se comparan con su signo, no en valor absoluto. La pendiente del IMg suele ser negativa, aunque con demandas muy convexas podría no serlo.} Se cumple siempre con CMg creciente; a diferencia de la empresa competitiva, el monopolista también puede maximizar en el tramo decreciente del CMg, siempre que este corte al IMg desde abajo, es decir, que decrezca más despacio que el IMg (anexo~\ref{anexo:cso}).\footnote{Si la función de beneficio es estrictamente cuasicóncava, la cantidad que maximiza el beneficio es única.}
\item \emph{No hay curva de oferta}. En competencia perfecta, $p=\mathrm{CMg}$ asocia a cada precio una única cantidad ofrecida. En el monopolio, la cantidad depende también de la forma de la demanda, a través de su elasticidad: 2 demandas distintas pueden dar el mismo precio con cantidades distintas, o la misma cantidad con precios distintos. No hay una relación única entre precio y cantidad, es decir, no hay curva de oferta propiamente dicha.
\item \emph{Condición de viabilidad}. Es la condición de participación de \eqref{eq:monopolio}, evaluada en el óptimo: además de la condición de primer orden, el monopolista solo produce si el beneficio de producir $q^M$ no es menor que el de cerrar:
  \begin{itemize}
  \item a \emph{corto plazo}, con algún factor fijo, hay costes fijos que se pagan produzca o no \vertema{3.A.12}. Produce $q^M$ si $p(q^M)\geq\mathrm{CVMe}(q^M)$: puede tener pérdidas, pero nunca mayores que el coste fijo;\footnote{Precisión sobre la condición de participación de \eqref{eq:monopolio}: $B(q)\geq 0$ es la de largo plazo; a corto plazo basta $p\geq\mathrm{CVMe}$, porque el coste fijo se paga igual, produzca o no, y lo que se compara es $B(q^M)\geq B(0)=-\mathrm{CF}$. El monopolista puede tener beneficios (gráfico de la izquierda) o pérdidas (derecha) a corto plazo; con pérdidas, sigue produciendo mientras el precio cubra el coste variable medio.\graficonota[0.6]{monopolio-corto-plazo}}
  \item a \emph{largo plazo}, todos los factores son variables y todos los costes evitables \vertema{3.A.12}: produce solo si $p(q^M)\geq\mathrm{CMe}(q^M)$, con beneficios ($p>\mathrm{CMe}$) o sin ellos ($p=\mathrm{CMe}$), pero nunca con pérdidas;\footnote{A largo plazo el monopolista tiene tiempo para ampliar su planta o para usar la existente en el nivel que maximiza su beneficio: puede corregir las pérdidas que sufra a corto plazo cambiando el tamaño de la planta.}
  \item la condición se evalúa \emph{en la cantidad óptima}, no en el mínimo del coste medio.\footnote{Las 2 formulaciones que se encuentran en los manuales coinciden en lo que importa. Que la demanda pase por encima del CMe en algún punto es necesario para que exista algún tramo rentable; pero, si existe un $q$ con $p(q)\geq\mathrm{CMe}(q)$, entonces $B(q)\geq 0$ y, como $q^M$ maximiza el beneficio, $B(q^M)\geq 0$, es decir, $p(q^M)\geq\mathrm{CMe}(q^M)$. Basta, por tanto, comprobarla en el óptimo. Lo que no es correcto es exigir que el precio de monopolio supere el \emph{mínimo} del CMe: el monopolista no produce en ese mínimo, sino donde $\mathrm{IMg}=\mathrm{CMg}$.} El anexo~\ref{anexo:viabilidad} lo muestra gráficamente.
  \end{itemize}
\item \emph{Interpretación gráfica} (gráfico~\ref{fig:equilibrio-monopolio}): el IMg corta al CMg en $q^M$, y el precio se lee en la demanda, $p^M=p(q^M)$, por encima del coste marginal. La cantidad competitiva, $q^c$, es la que iguala la demanda y el coste marginal, con precio $p^c$.
\end{itemize}

\begin{figure}[H]
\centering
\caption{Equilibrio del monopolista: $\mathrm{IMg}=\mathrm{CMg}$ y precio sobre la demanda}\label{fig:equilibrio-monopolio}
\grafico[0.6\textwidth]{equilibrio-monopolio}
\fuente{elaboración propia a partir de \textcite[cap.~1]{Tirole1988Theory}.}
\end{figure}

\subparagraph{iii) Implicaciones}

\begin{itemize}
\item En el equilibrio, el monopolista produce hasta igualar el ingreso marginal y el coste marginal, y fija el precio según la demanda: $p^M>\mathrm{CMg}(q^M)$ y $q^M<q^c$. Cobra más y produce menos que la competencia.
\item A diferencia de la competencia perfecta:
  \begin{enumerate}[label=\alph*)]
  \item puede tener \emph{beneficios extraordinarios a largo plazo}: atraerían competidores, pero las barreras de entrada les impiden entrar. Dependen de las barreras y de que la demanda no sea perfectamente elástica (es la lectura económica del índice de \autor{Lerner});
  \item \emph{no produce necesariamente en el mínimo del coste medio a largo plazo}: según el tamaño de la demanda, puede quedarse por debajo de la escala eficiente, con capacidad ociosa y economías de escala sin aprovechar, o por encima, con deseconomías.\footnote{En el gráfico de la izquierda, $\mathrm{IMg}=\mathrm{CMg}$ se cumple justo en el mínimo del CMe, con pleno aprovechamiento de la capacidad y de las economías de escala. En el del centro, en el tramo decreciente del CMe: un mercado pequeño, con exceso de capacidad. En el de la derecha, en el tramo creciente: un mercado muy grande que obliga a producir más allá de la escala óptima, con deseconomías de escala.\graficonota[0.85]{escala-monopolio}}
  \end{enumerate}
\end{itemize}

\paragraph{c) Extensiones}

\begin{itemize}
\item Una aplicación directa del modelo es el efecto de los impuestos sobre el monopolista, que el test suele preguntar (anexo~\ref{anexo:test}).
\end{itemize}

\begin{anotaciones}
\textbf{Impuestos sobre el monopolista}. Con un impuesto unitario $t$, el problema es $\max_q\ p(q)\,q-C(q)-t\,q$ y la condición de primer orden pasa a $\mathrm{IMg}(q)=\mathrm{CMg}(q)+t$: la cantidad baja y el precio sube. Cuánto sube (la \emph{traslación} del impuesto) depende de la curvatura de la demanda, no solo de su pendiente \parencite{BulowPfleiderer1983Note}: con demanda lineal $p=a-bq$ y coste marginal constante $c$, $p^M=(a+c+t)/2$, y el precio sube la mitad del impuesto; con demanda isoelástica, $p^M=\frac{|\varepsilon_{q,p}|}{|\varepsilon_{q,p}|-1}\,(c+t)$, y sube \emph{más} que el impuesto. Con un impuesto \emph{ad valorem} $\tau$ sobre el precio, la condición es $(1-\tau)\,\mathrm{IMg}(q)=\mathrm{CMg}(q)$ y también reduce la cantidad. Un impuesto sobre el beneficio (con un tipo menor que 1) o de suma fija no altera la cantidad, porque no cambia la condición de primer orden.
\end{anotaciones}

\paragraph{d) Valoración}

\begin{itemize}
\item El modelo explica con una sola condición, $\mathrm{IMg}=\mathrm{CMg}$, por qué el monopolista cobra más y produce menos, y da una medida del poder de mercado, el índice de \autor{Lerner}, que se usa en la práctica. Ha fundamentado la intervención en los mercados monopolísticos y la eliminación de barreras legales a la entrada.
\item Pero deja abierta la pregunta normativa: ¿cuánto le cuesta a la sociedad que el precio supere el coste marginal? Es lo que analizamos a continuación.
\end{itemize}

% --------------------------------------------------------------------------
\subsection{Análisis de eficiencia y bienestar}\label{sec:analisis-eficiencia-bienestar}

\begin{itemize}
\item Una vez visto cómo fija el monopolista el precio y la cantidad, pasamos a preguntarnos cuánto le cuesta eso a la sociedad. La \emph{conclusión anticipada}: el coste estático, el triángulo de \autor{Harberger}, es pequeño en las estimaciones clásicas, pero el coste total puede ser mucho mayor si se suman los costes inflados y los recursos que se gastan en conseguir y defender el monopolio.
\end{itemize}

\paragraph{a) Idea}

\begin{itemize}
\item Comparamos el equilibrio del monopolio con el competitivo, que maximiza el \concepto{excedente total}, $\mathrm{ET}=\mathrm{EC}+\mathrm{EP}$: es el 1TFEB en equilibrio parcial, que exige $p=\mathrm{CMg}$ en la última unidad intercambiada \vertema{3.A.16} \vertema{3.A.22}.\footnote{En equilibrio parcial, con utilidad cuasilineal, una asignación es eficiente en el sentido de \autor{Pareto} si y solo si maximiza el excedente total, y eso exige $p=\mathrm{CMg}$ en la última unidad (\emph{eficiencia asignativa}): aquí las 2 nociones coinciden. Quién se queda con el excedente es una cuestión de equidad, no de eficiencia \vertema{3.A.16}.}
\end{itemize}

\paragraph{b) Modelo}

\subparagraph{i) Supuestos}

\begin{itemize}
\item La demanda $p(q)$ mide la valoración marginal de los consumidores: con utilidad cuasilineal, el excedente del consumidor es una medida exacta de su bienestar \vertema{3.A.9}.
\item El coste marginal mide el coste de oportunidad de producir.
\item El monopolio y la industria competitiva tienen la misma función de costes (supuesto que se levanta con la ineficiencia X).
\end{itemize}

\subparagraph{ii) Desarrollo}

\begin{itemize}
\item \emph{Ineficiencia asignativa en el mercado del bien} (gráfico~\ref{fig:bienestar-monopolio}). Al pasar de la competencia al monopolio:
  \begin{itemize}
  \item el rectángulo entre $p^c$ y $p^M$ hasta $q^M$ pasa de los consumidores al monopolista: es una \emph{transferencia}, no una pérdida;
  \item los 2 triángulos correspondientes a las unidades entre $q^M$ y $q^c$, uno que era excedente del consumidor y otro que era excedente del productor, desaparecen: son la \concepto{pérdida irrecuperable de eficiencia} o \concepto{triángulo de \autor{Harberger}}. Son unidades que los consumidores valoran por encima de lo que cuesta producirlas y que no se producen.
  \end{itemize}
\item \emph{Cuánto vale el triángulo}. Con coste marginal constante, aproximándolo por un triángulo con la elasticidad del equilibrio,
  \begin{equation}
  \mathrm{PIE}\approx\tfrac12\,\Delta p\,\Delta q\approx\tfrac12\,|\varepsilon_{q,p}|\,\mathcal{L}^2\,p^Mq^M, \label{eq:harberger}
  \end{equation}
  proporcional al \emph{cuadrado} del margen. Y como el monopolista fija su precio según \eqref{eq:lerner}, $|\varepsilon_{q,p}|=1/\mathcal{L}$ y $\mathrm{PIE}\approx\tfrac12\,\mathcal{L}\,p^Mq^M=\tfrac12\,B$: la pérdida irrecuperable es del orden de la mitad del beneficio del monopolista (exactamente la mitad con demanda lineal, coste marginal constante y sin costes fijos).\footnote{Con $p=a-bq$ y $\mathrm{CMg}=c$: $q^M=(a-c)/2b$, $p^M=(a+c)/2$, $B=(a-c)^2/4b$ y $\mathrm{PIE}=\frac{1}{2}\cdot\frac{a-c}{2}\cdot\frac{a-c}{2b}=(a-c)^2/8b=B/2$.}
\end{itemize}

\begin{figure}[H]
\centering
\caption{El coste social del monopolio: transferencia, triángulo de \autor{Harberger}, búsqueda de rentas e ineficiencia X}\label{fig:bienestar-monopolio}
\grafico[0.62\textwidth]{bienestar-monopolio}
\fuente{elaboración propia a partir de \textcite{Harberger1954Monopoly}, \textcite{Posner1975Social} y \textcite{Leibenstein1966Allocative}.}
\end{figure}

\begin{itemize}
\item \emph{Otras ineficiencias}. El triángulo es solo la parte más visible del coste social del monopolio; hay otras 4, que pueden ser mayores:
  \begin{itemize}
  \item \emph{Ineficiencia asignativa en los mercados de factores}: el monopolista contrata cada factor hasta que su \concepto{ingreso del producto marginal}, $\mathrm{IPMg}_i=\mathrm{IMg}\cdot\mathrm{PMg}_i$, iguala su precio, y no hasta que lo iguala el \concepto{valor del producto marginal}, $\mathrm{VPMg}_i=p\cdot\mathrm{PMg}_i$, que es mayor porque $\mathrm{IMg}<p$. Emplea menos factores de los eficientes: es lo que \autor{Robinson} llama \emph{explotación monopolística} de los factores.\footnote{La combinación de factores sigue siendo la más barata para la cantidad producida: $\mathrm{RMST}_{LK}=w_L/w_K$. Aquí suponemos que el monopolista compra los factores en mercados competitivos, junto con otros demandantes; si además es el único comprador (monopsonista), el precio del factor depende de la cantidad que compra (\ref{sec:monopsonio}). Tampoco se dice que la retribución de los factores sea menor que en competencia, porque con menos producción su productividad marginal es mayor, sino que es menor que el valor de su producto marginal.}
  \item \emph{Ineficiencia productiva}: no tiene por qué producir en el mínimo del coste medio a largo plazo (\ref{sec:equilibrio-monopolio}), así que puede operar con una escala ineficiente.
  \item \emph{Ineficiencia X} \parencite{Leibenstein1966Allocative}: sin la presión de la competencia, los costes observados se alejan del mínimo, $C^{\mathrm{obs}}=(1+\chi)\,C(q)$ con $\chi\geq 0$ \vertema{3.A.12}. Como dice \textcite[8]{Hicks1935Annual}, «el mejor de todos los beneficios del monopolio es una vida tranquila». En el gráfico~\ref{fig:bienestar-monopolio}, el coste marginal sube a $(1+\chi)\,\mathrm{CMg}$ y se pierde, además, la franja entre las 2 curvas de coste hasta $q^M$: recursos que se gastan en producir lo que podría producirse con menos.
  \item \emph{Búsqueda de rentas}: el beneficio del monopolio es un premio por el que se compite (grupos de presión, licencias, litigios). \textcite{Tullock1967Welfare} señala que los recursos dedicados a conseguirlo son un coste social, y \textcite{Posner1975Social} supone que la competencia por el premio lo disipa entero: todo el rectángulo del beneficio se suma a la pérdida \vertema{4.B.2}.
  \end{itemize}
\end{itemize}

\subparagraph{iii) Implicaciones}

\begin{itemize}
\item La pérdida del monopolio va del triángulo de \autor{Harberger} (cota inferior) al triángulo más la franja de la ineficiencia X y el rectángulo de la búsqueda de rentas (cota superior).
\item Como la pérdida crece con el cuadrado del margen, los monopolios con demanda muy inelástica, que cargan márgenes altos, son los más costosos.
\end{itemize}

\paragraph{c) Extensiones: eficiencia dinámica}

\begin{itemize}
\item \textcite{Schumpeter1942Capitalism} defiende que el beneficio del monopolio financia y premia la innovación: la pérdida estática se compensa con más progreso técnico. De ahí las patentes, un monopolio temporal concedido a propósito.
\item \textcite{Arrow1962Economic} responde que el incentivo a innovar es mayor en competencia, porque el monopolista ya cobra rentas con la tecnología vieja y solo gana la diferencia (\emph{efecto reemplazo}).
\item La evidencia sugiere una relación en forma de U invertida entre competencia e innovación: la innovación aumenta con la competencia cuando esta es escasa y disminuye cuando es muy intensa \parencite{AghionEtAl2005Competition} \vertema{3.A.44}.
\end{itemize}

\paragraph{d) Evidencia empírica}

\begin{itemize}
\item \textcite{Harberger1954Monopoly}, con datos de la industria manufacturera de Estados Unidos de 1924-1928 y elasticidad unitaria, estima que la pérdida no llegaba a «una treceava parte del uno por ciento de la renta nacional» \parencite{Schmitz2012New}. \textcite{Posner1975Social}, al añadir la búsqueda de rentas, y \textcite{CowlingMueller1978Social}, con la fórmula $\mathrm{PIE}=B/2$ y los gastos en publicidad, obtienen pérdidas muy superiores.
\item Hoy los márgenes son mayores: \textcite{DeLoeckerEeckhoutUnger2020Rise} estiman que el margen medio de las empresas cotizadas de Estados Unidos sobre el coste marginal pasó del 21~\% en 1980 al 61~\% en 2016, sobre todo porque las empresas con más poder de mercado ganaron cuota. Es un margen agregado de la economía, no de monopolios, pero es justo el índice de \autor{Lerner} que mide el poder de mercado.
\end{itemize}

\paragraph{e) Valoración}

\begin{itemize}
\item El monopolio suele ser ineficiente, pero de ahí no se sigue que haya que desmantelarlo siempre: hay monopolios eficientes en la asignación, como el que discrimina perfectamente (\ref{sec:discriminacion-precios}), o en los costes, como el monopolio natural (\ref{sec:monopolio-natural}), y la intervención pública tiene sus propios costes \vertema{3.A.23}. Lo que haya que hacer depende del tipo de monopolio y del origen de su poder.
\item \emph{Recapitulando}: $\mathrm{IMg}=\mathrm{CMg}$, un margen inversamente proporcional a la elasticidad y una pérdida de bienestar que va más allá del triángulo de \autor{Harberger}. Pero el modelo supone un solo precio, unos costes sin ventajas de tamaño y un solo producto. Pasamos a levantar esos supuestos.
\end{itemize}


% ==========================================================================
\section{Desarrollos de la teoría del monopolio}\label{sec:desarrollos-teoria-monopolio}
\minutos{11,5}

\begin{notaopositor}
En el cante: la definición de discriminación, sus 3 condiciones y la clasificación por la información del monopolista; el primer grado como una capa más del gráfico del equilibrio; el segundo grado con el modelo de 2 tipos y su gráfico (restricciones saturadas, renta informativa, «no distorsión arriba», distorsión abajo) y la calidad en una frase; la tarifa en dos partes en una frase, señalando en su gráfico lo que se apropia el monopolista y lo que no; el tercer grado con la regla de \autor{Lerner} por grupos y su juicio de bienestar; la conjetura de \autor{Coase} en una frase (\tiempo{5,5}). El monopolio natural: la subaditividad como definición, el gráfico de la sostenibilidad, las economías de alcance como condición necesaria y los subsidios cruzados; la regulación: el dilema, el primer óptimo con pérdidas, el coste medio, la tarifa en dos partes y la regla de \autor{Ramsey}-\autor{Boiteux} junto al índice de \autor{Lerner}, con su lectura económica (\tiempo{4,5}). La producción conjunta: la regla de \autor{Lerner} con elasticidades cruzadas (solo su estructura y los 2 signos), las proporciones fijas en una frase ($\mathrm{IMg}_1+\mathrm{IMg}_2=\mathrm{CMg}$ y el desecho) y el multiplanta (\tiempo{1,5}). Los 3 casos de la tarifa en dos partes en detalle, la venta conjunta, los ejemplos, la sostenibilidad con reacción del incumbente, el aprendizaje y el modelo completo de las proporciones fijas (anexo~\ref{anexo:proporciones-fijas}) son de profundización; las derivaciones, en el anexo~\ref{anexo:derivaciones}.
\end{notaopositor}

\begin{itemize}
\item Una vez estudiado el modelo básico, pasamos a sus desarrollos. La idea que los une, y \emph{conclusión anticipada} del bloque: cada vez que se levanta un supuesto, el monopolista adapta la misma regla de la elasticidad inversa, y el juicio de bienestar puede cambiar de signo: la discriminación puede aumentar la producción, y un monopolio natural puede ser la forma más barata de producir.
\end{itemize}

% --------------------------------------------------------------------------
\subsection{Discriminación de precios}\label{sec:discriminacion-precios}

\begin{itemize}
\item En el modelo básico el monopolista cobra un precio uniforme. Pero puede cobrar precios distintos para extraer parte del excedente que el precio uniforme deja a los consumidores. Hay \concepto{discriminación de precios} cuando se venden unidades del mismo bien a precios distintos, al mismo o a distintos compradores, sin que la diferencia se deba a diferencias de coste \parencite{Varian1989Price}.
\item Requiere 3 condiciones:
  \begin{enumerate}
  \item \emph{poder de mercado}: con demanda perfectamente elástica no hay excedente que extraer;
  \item \emph{información} sobre la disposición a pagar de los consumidores, directa, a través de una señal o a través de lo que eligen;
  \item \emph{ausencia de reventa} (arbitraje) entre consumidores: si quien compra barato pudiera revender a quien paga caro, solo quedaría un precio.
  \end{enumerate}
\item \textcite{Pigou1920Economics} distingue 3 grados: en el primero, cada unidad se vende a su precio de demanda; en el segundo, hay precios por tramos de cantidad; en el tercero, precios distintos para grupos separados por alguna señal observable. Hoy se ordenan por la información del monopolista \parencite[cap.~3]{Tirole1988Theory}:
  \begin{itemize}
  \item \concepto{primer grado} (discriminación perfecta): observa la valoración de cada consumidor;
  \item \concepto{segundo grado} (fijación no lineal de precios): no observa la valoración de cada consumidor, pero conoce su distribución, y ofrece a todos el mismo menú, del que cada uno elige: los consumidores se \emph{autoseleccionan};
  \item \concepto{tercer grado} (por grupos): observa una característica del consumidor (la edad, el lugar de residencia) correlacionada con la valoración, y cobra a cada grupo un precio lineal distinto.
  \end{itemize}
\item La discriminación \emph{intertemporal} de los bienes duraderos, que algunos manuales llaman «de cuarto grado», no es un grado más de la clasificación de \autor{Pigou}: es una forma de autoselección en el tiempo (\ref{sec:discriminacion-intertemporal-bienes}).
\end{itemize}

\subsubsection{Discriminación de primer grado}\label{sec:discriminacion-primer-grado}

\paragraph{a) Idea}

\begin{itemize}
\item Es el caso extremo: el monopolista sabe cuánto está dispuesto a pagar cada consumidor por cada unidad y puede cobrárselo.
\end{itemize}

\paragraph{b) Modelo}

\subparagraph{i) Supuestos}

\begin{itemize}
\item El monopolista observa la valoración de cada consumidor por cada unidad y puede cobrar por ella exactamente su precio de reserva; no hay reventa.
\end{itemize}

\subparagraph{ii) Desarrollo}

\begin{itemize}
\item Cada unidad adicional se vende a su precio de reserva \emph{sin bajar el de las demás}, así que el ingreso marginal es la propia demanda, $\mathrm{IMg}(q)=p(q)$, y produce hasta que la última unidad se vende a su coste marginal, $p(q^{*})=\mathrm{CMg}(q^{*})$: la cantidad competitiva, $q^{*}=q^c$ (gráfico~\ref{fig:primer-grado}).
\end{itemize}

\begin{figure}[H]
\centering
\caption{Discriminación de primer grado: la demanda es el ingreso marginal y todo el excedente es del monopolista}\label{fig:primer-grado}
\grafico[0.55\textwidth]{primer-grado}
\fuente{elaboración propia a partir de \textcite[cap.~3]{Tirole1988Theory}.}
\end{figure}

\subparagraph{iii) Implicaciones}

\begin{itemize}
\item La asignación es \emph{eficiente}: el excedente total es el máximo y el triángulo de \autor{Harberger} desaparece. Pero todo es excedente del productor: el excedente del consumidor es nulo. El problema es de \emph{equidad}, no de eficiencia, y la intervención solo se justifica por motivos distributivos.\footnote{La última unidad se vende a $p=\mathrm{CMg}$, así que la asignación es eficiente: en equilibrio parcial, la eficiencia asignativa y la de \autor{Pareto} coinciden aquí. El problema es solo distributivo, porque los consumidores se quedan sin excedente.}
\item Es la cota superior del beneficio de cualquier forma de discriminación.
\end{itemize}

\paragraph{c) Evidencia empírica}

\begin{itemize}
\item Es difícil de encontrar en estado puro: el regateo en un zoco o el médico rural que cobra a cada paciente según lo que puede pagar se le acercan. Hoy, la \emph{tarificación personalizada} con datos y algoritmos se aproxima a ella, aunque la evidencia sobre su uso masivo es todavía limitada.
\end{itemize}

\paragraph{d) Valoración}

\begin{itemize}
\item Exige una información que el monopolista casi nunca tiene. Cuando no observa la valoración de cada uno, tiene que conseguir que los consumidores la revelen: es la discriminación de segundo grado.
\end{itemize}

\subsubsection{Discriminación de segundo grado: autoselección en cantidad y en calidad}\label{sec:discriminacion-segundo-grado}

\paragraph{a) Idea}

\begin{itemize}
\item Ahora el monopolista no sabe quién es quién, y si cobrara más a los que más valoran el bien, estos se harían pasar por los que menos lo valoran. Lo que sí puede hacer es ofrecer a todos el mismo \emph{menú} de opciones (paquetes de cantidad y pago, o versiones de calidad y precio), diseñado para que cada tipo elija la suya y revele así su valoración. Es un problema de \concepto{autoselección} (\emph{screening}) en un entorno de selección adversa \vertema{3.A.13} \parencites{MussaRosen1978Monopoly}{MaskinRiley1984Monopoly}.
\item Las personas que compran la misma cantidad pagan lo mismo, pero el precio por unidad cambia con la cantidad: es una \concepto{tarifa no lineal}.
\end{itemize}

\paragraph{b) Modelo}

\subparagraph{i) Supuestos}

\begin{itemize}
\item 2 tipos de consumidores, con valoración baja $t_1$ y alta $t_2>t_1$, en proporciones $\pi$ y $1-\pi$, que el monopolista conoce; no sabe de qué tipo es cada uno.
\item El tipo $i$ obtiene una utilidad $t_i\,u(q)-T$ al comprar $q$ unidades pagando $T$ en total, con $u(0)=0$, $u'>0$ y $u''<0$; si no compra, 0.
\item Coste marginal constante $c$. El monopolista ofrece un menú $\{(q_1,T_1),(q_2,T_2)\}$ y no hay reventa.
\end{itemize}

\subparagraph{ii) Desarrollo}

\begin{itemize}
\item El monopolista elige el menú que maximiza su beneficio esperado:
  \begin{equation}
  \optim{\max}{q_1,T_1,q_2,T_2}{\pi\,(T_1-c\,q_1)+(1-\pi)\,(T_2-c\,q_2)}{t_1u(q_1)-T_1\geq 0\quad(\mathrm{PC}_1)\\ t_2u(q_2)-T_2\geq 0\quad(\mathrm{PC}_2)\\ t_1u(q_1)-T_1\geq t_1u(q_2)-T_2\quad(\mathrm{IC}_1)\\ t_2u(q_2)-T_2\geq t_2u(q_1)-T_1\quad(\mathrm{IC}_2)} \label{eq:segundo-grado}
  \end{equation}
  con 2 \emph{restricciones de participación} (PC: cada tipo prefiere comprar a no comprar) y 2 \emph{de compatibilidad de incentivos} (IC: cada tipo prefiere su opción a la del otro).
\item En el óptimo solo se saturan 2 (anexo~\ref{anexo:segundo-grado}): la de participación del tipo bajo y la de incentivos del alto. Al tipo bajo se le extrae todo su excedente, $T_1=t_1u(q_1)$; al alto hay que dejarle lo que ganaría haciéndose pasar por bajo, su \concepto{renta informativa}:
  \begin{equation}
  T_2=t_2u(q_2)-\underbrace{(t_2-t_1)\,u(q_1)}_{\text{renta informativa del tipo alto}} .
  \end{equation}
\item Sustituyendo los pagos y derivando respecto a las cantidades:
  \begin{equation}
  t_2\,u'(q_2)=c,\qquad t_1\,u'(q_1)\left[1-\frac{1-\pi}{\pi}\cdot\frac{t_2-t_1}{t_1}\right]=c\ \Longrightarrow\ t_1\,u'(q_1)>c . \label{eq:autoseleccion}
  \end{equation}
\item \emph{Interpretación gráfica} (gráfico~\ref{fig:autoseleccion}): con las demandas de los 2 tipos, $t_iu'(q)$, y el coste marginal $c$, el tipo alto compra donde su demanda corta a $c$; el tipo bajo, a la izquierda de ese punto; la renta informativa es el área entre las 2 demandas hasta $q_1$.
\end{itemize}

\begin{figure}[H]
\centering
\caption{Discriminación de segundo grado con 2 tipos: cantidad eficiente para el tipo alto, distorsión abajo y renta informativa}\label{fig:autoseleccion}
\grafico[0.6\textwidth]{autoseleccion}
\fuente{elaboración propia a partir de \textcite{MaskinRiley1984Monopoly} y \textcite[cap.~3]{Tirole1988Theory}.}
\end{figure}

\subparagraph{iii) Implicaciones}

\begin{itemize}
\item \emph{No hay distorsión arriba}: el tipo alto compra la cantidad eficiente (su valoración marginal iguala el coste marginal) y se queda con una renta informativa.
\item \emph{Distorsión abajo}: al tipo bajo se le ofrece \emph{menos} de lo eficiente. Cada unidad que se le quita le cuesta al monopolista algo de excedente que ya no puede cobrarle, pero reduce la renta que tiene que dejar al tipo alto, porque la opción del bajo se vuelve menos atractiva de imitar. Cuantos más tipos altos haya ($(1-\pi)/\pi$ grande) y más distintos sean, mayor es la distorsión; si $\pi\,t_1\leq(1-\pi)\,(t_2-t_1)$, conviene excluir al tipo bajo.
\item \emph{Presentación del menú}: equivale a una tarifa no lineal con \emph{descuentos por cantidad} (el paquete grande tiene un precio medio menor), o como 2 tarifas en dos partes entre las que se elige: cuota mayor y precio por unidad igual a $c$ para el tipo alto; cuota menor y precio por unidad mayor que $c$ para el bajo.
\item \emph{Bienestar}: la cantidad del tipo alto es eficiente y la del bajo no; frente al precio uniforme, la autoselección puede aumentar o reducir el excedente total.
\end{itemize}

\paragraph{c) Extensiones}

\begin{itemize}
\item \emph{Discriminación en calidad}. Si en lugar de cantidades se ofrecen calidades con un coste creciente, el análisis es idéntico: el tipo alto recibe la calidad eficiente y el bajo una calidad \emph{degradada} a propósito, para que la versión barata no tiente al alto \parencite{MussaRosen1978Monopoly}. Ejemplos: la tapa dura y la tapa blanda de un libro; la clase \emph{business} y la turista; las versiones «básica» y «profesional» de un programa, a veces fabricando primero la completa y quitándole funciones después (\emph{damaged goods}, \textcite{DeneckereMcAfee1996Damaged}).\footnote{Xavier \autor{Sala-i-Martín} llama \emph{paradoja de Barbie} a que muñecas que cuestan lo mismo de fabricar se vendan a precios muy distintos según la profesión que representan \parencite{SalaiMartin2013Barbie}. Más que discriminación de tercer grado, es un ejemplo de versiones entre las que el comprador elige.}
\item \emph{Una sola tarifa en dos partes para todos}, $T(q)=A+p\,q$: una cuota fija $A$ que da derecho a comprar y un precio $p$ por unidad, de modo que el precio medio baja con la cantidad.\footnote{El ejemplo clásico es el parque de atracciones, que cobra una entrada y un precio por atracción \parencite{Oi1971Disneyland}: el dilema es cobrar una entrada cara y atracciones baratas o al revés.} El gráfico~\ref{fig:tarifa-dos-partes} muestra 3 casos con 2 tipos de consumidores:
  \begin{enumerate}[label=\roman*)]
  \item si el monopolista puede distinguir a los tipos y cobrarles cuotas distintas, fija $p=c$ y cuotas iguales al excedente de cada tipo, $A_i=\mathrm{EC}_i(c)$: equivale a la discriminación perfecta;
  \item si la cuota tiene que ser la misma para todos y no quiere excluir a nadie, la fija en el excedente del tipo bajo, $A=\mathrm{EC}_1(p)$, y el precio por unidad queda \emph{por encima} del coste marginal \parencite{Oi1971Disneyland}. Con $n_i$ consumidores de cada tipo y demandas $q_i(p)$,
    \begin{equation}
    p-c=\frac{n_2\,[q_2(p)-q_1(p)]}{-[n_1q_1'(p)+n_2q_2'(p)]}>0 :
    \end{equation}
    subir $p$ por encima de $c$ resta algo de excedente al tipo bajo, y por tanto a la cuota, pero permite cobrar al alto, por unidad, parte de lo que no se le puede cobrar en la cuota;
  \item si los tipos bajos son pocos, puede compensar fijar la cuota en el excedente del tipo alto, $A=\mathrm{EC}_2(c)$ con $p=c$, y excluir a los bajos.
  \end{enumerate}
  En el gráfico se ve también lo que el monopolista \emph{no} se apropia: nada en el caso~i); en el ii), la renta del tipo alto (lo que valora por encima de la cuota común y de lo que paga por unidad) y una pérdida irrecuperable en los 2 tipos, porque $p>c$; en el iii), el excedente que habría generado el tipo bajo, que se pierde al quedar excluido.
\end{itemize}

\begin{figure}[H]
\centering
\caption{Tarifa en dos partes con 2 tipos: con cuotas distintas, con una cuota común que no excluye y con una cuota que excluye al tipo bajo. En dorado, lo que se apropia el monopolista; en morado, lo que se queda el consumidor, y en rojo, la pérdida irrecuperable}\label{fig:tarifa-dos-partes}
\grafico[0.95\textwidth]{tarifa-dos-partes}
\fuente{elaboración propia a partir de \textcite{Oi1971Disneyland}.}
\end{figure}

\begin{anotaciones}
\textbf{Venta conjunta (\emph{bundling})}. Si el monopolista vende varios bienes, puede ofrecerlos en un paquete a un precio menor que la suma de los precios por separado. No es discriminación de segundo grado por cantidad, sino otra forma de discriminar: cuando las valoraciones de los consumidores están correlacionadas negativamente, el paquete reduce la dispersión de la disposición a pagar y permite extraer más excedente \parencite{AdamsYellen1976Commodity}. Ejemplo con coste nulo: el consumidor 1 valora el bien A en 6 y el B en 4; el 2, al revés. Por separado, lo mejor es cobrar 4 por cada bien (ingreso 16); el paquete a 10 lo compran los 2 (ingreso 20). En la \emph{venta conjunta mixta} se ofrecen el paquete y los bienes sueltos, y los consumidores se autoseleccionan, como en el segundo grado. Está muy extendida en las tecnologías de la información, por ejemplo en las suscripciones a revistas en línea. Si lo hace una empresa dominante para extender su poder a otro mercado, puede ser un abuso (ventas vinculadas) \vertema{3.B.39}.
\end{anotaciones}

\paragraph{d) Evidencia empírica}

\begin{itemize}
\item Las tarifas de telefonía móvil y de datos, los bonos de transporte y las suscripciones con varios niveles son menús de autoselección; \textcite{Varian1989Price} repasa los estudios empíricos.
\end{itemize}

\paragraph{e) Valoración}

\begin{itemize}
\item Explica prácticas generalizadas, como los descuentos por volumen o las versiones de un producto, que no se entienden con un precio uniforme, y muestra el coste de la información privada: la renta informativa y la distorsión abajo. Es la misma lógica que el menú de contratos del regulador (anexo~\ref{anexo:regulacion}). Su límite es que el monopolista debe conocer la distribución de los tipos; cuando lo que observa es una característica del consumidor, pasamos a la discriminación de tercer grado.
\end{itemize}

\subsubsection{Discriminación de tercer grado}\label{sec:discriminacion-tercer-grado}

\paragraph{a) Idea}

\begin{itemize}
\item El monopolista no conoce la valoración de cada consumidor, pero observa una característica suya (edad, lugar de residencia, condición de estudiante) que lo sitúa en un grupo, y conoce la demanda de cada grupo. Cobra un precio distinto a cada grupo, el mismo para todas las unidades dentro de él.
\item Además de la ausencia de reventa entre grupos, necesita poder separar la demanda en varios submercados y que sus elasticidades sean distintas.\footnote{Es una situación de información asimétrica \vertema{3.A.13}: hay información oculta (la valoración de cada consumidor), heterogeneidad (si todos fueran iguales, no habría nada que discriminar) y conflicto de intereses (el consumidor no tiene interés en revelar que valora mucho el bien). La característica observable es una señal imperfecta de la valoración.}
\end{itemize}

\paragraph{b) Modelo}

\subparagraph{i) Supuestos}

\begin{itemize}
\item 2 grupos, $k=1,2$, con demandas inversas $p_k(q_k)$ conocidas; precio lineal dentro de cada grupo; sin reventa entre grupos; un único coste, $C(q_1+q_2)$.
\end{itemize}

\subparagraph{ii) Desarrollo}

\begin{itemize}
\item Es un monopolio que vende en 2 mercados separados:
  \begin{equation}
  \optim{\max}{q_1,q_2}{p_1(q_1)\,q_1+p_2(q_2)\,q_2-C(q_1+q_2)}{q_1\geq 0\\ q_2\geq 0}
  \end{equation}
\item Las condiciones de primer orden igualan el ingreso marginal de cada mercado con el coste marginal común, y de ahí sale una regla de \autor{Lerner} por grupo:
  \begin{equation}
  \mathrm{IMg}_1(q_1)=\mathrm{IMg}_2(q_2)=\mathrm{CMg}(q_1+q_2)\ \Longrightarrow\ \mathcal{L}_k=\frac{p_k-\mathrm{CMg}}{p_k}=\frac{1}{|\varepsilon_{q_k,p_k}|},\quad k=1,2. \label{eq:tercer-grado}
  \end{equation}
  Si en un mercado el IMg fuera mayor que en el otro, le convendría pasar una unidad de este a aquel.
\item \emph{Interpretación gráfica} (gráfico~\ref{fig:tercer-grado}): el IMg agregado (suma horizontal de los IMg de los 2 grupos) corta al CMg y fija la producción total; la altura de ese corte da en cada mercado la cantidad y, en su demanda, el precio.
\end{itemize}

\begin{figure}[H]
\centering
\caption{Discriminación de tercer grado: precio más alto en el mercado de demanda menos elástica}\label{fig:tercer-grado}
\grafico[0.95\textwidth]{tercer-grado}
\fuente{elaboración propia a partir de \textcite{Robinson1933Economics}.}
\end{figure}

\subparagraph{iii) Implicaciones}

\begin{itemize}
\item Es la regla de la elasticidad inversa aplicada a cada grupo: \emph{el precio es mayor en el grupo de demanda menos elástica}. Si las elasticidades son iguales, no compensa discriminar.
\item \emph{Bienestar}. El beneficio no puede bajar, porque el precio uniforme sigue siendo posible; el excedente total puede subir o bajar. La discriminación desplaza ventas del grupo que paga más al que paga menos, lo que es ineficiente (la misma unidad se valora distinto en cada grupo), y por eso \emph{solo puede aumentar el excedente total si aumenta la cantidad total} \parencites{Schmalensee1981Output}{Varian1985Price}. Con coste marginal constante $c$ y $p^u$ el precio uniforme,
  \begin{equation}
  \textstyle\sum_k(p_k-c)\,\Delta q_k\ \leq\ \Delta\mathrm{ET}\ \leq\ (p^u-c)\,\Delta Q ,
  \end{equation}
  con $\Delta Q=\sum_k\Delta q_k$. Con demandas lineales y todos los grupos servidos, la cantidad total no cambia y el excedente total baja \parencite{Robinson1933Economics}.
\item El caso favorable es el de un grupo que, con precio uniforme, no compraría (el precio se fija para el grupo de valoración alta) y con discriminación sí: la cantidad aumenta y pueden ganar todos, una mejora en el sentido de \autor{Pareto}. Es el caso de los medicamentos más baratos en los países pobres.
\end{itemize}

\paragraph{c) Evidencia empírica}

\begin{itemize}
\item Los descuentos para jóvenes y mayores en el transporte, las tarifas pensadas para quienes viajan por negocios, los precios distintos por país de los medicamentos o de los libros de texto, y los cupones de descuento, que separan a los consumidores más sensibles al precio: los que se toman la molestia de usarlos.
\end{itemize}

\paragraph{d) Valoración}

\begin{itemize}
\item Es la forma de discriminación más frecuente y la más fácil de observar. Su efecto sobre el bienestar es ambiguo: la pregunta clave es si la discriminación abre mercados que con un precio uniforme no se atenderían.
\end{itemize}

\subsubsection{Discriminación intertemporal: bienes duraderos y conjetura de \autor{Coase}}\label{sec:discriminacion-intertemporal-bienes}

\paragraph{a) Idea}

\begin{itemize}
\item Si el bien dura y el monopolista puede bajar el precio con el tiempo, vende primero caro a los que más lo valoran y después más barato al resto: los consumidores se autoseleccionan por el momento de compra. Pero quien compra hoy no vuelve a comprar mañana, así que los bienes que vende en periodos distintos son sustitutivos: el monopolista compite con su yo futuro.
\end{itemize}

\paragraph{b) Modelo}

\subparagraph{i) Supuestos}

\begin{itemize}
\item El bien es duradero: quien compra hoy no vuelve a comprar mañana.
\item El monopolista puede revisar el precio con frecuencia y no puede comprometerse a mantenerlo.
\item Los consumidores son pacientes y anticipan los precios futuros.
\end{itemize}

\subparagraph{ii) Desarrollo}

\begin{itemize}
\item \concepto{Conjetura de \autor{Coase}}: \textcite{Coase1972Durability} sostiene que, si puede revisar el precio con frecuencia y los consumidores son pacientes, nadie paga hoy un precio alto sabiendo que mañana bajará, y el precio se acerca al coste marginal: el monopolista pierde su poder de mercado. \textcite{GulSonnenscheinWilson1986Foundations} lo demuestran con horizonte infinito y revisiones del precio cada vez más frecuentes.\footnote{El problema es de inconsistencia temporal: al monopolista le gustaría comprometerse a no bajar el precio, pero, una vez que ha vendido a los que más lo valoran, le conviene bajarlo para vender al resto, y los consumidores lo anticipan \vertema{3.A.14}.}
\end{itemize}

\subparagraph{iii) Implicaciones}

\begin{itemize}
\item El monopolista gana si consigue comprometerse a no bajar el precio. Por eso alquila en lugar de vender (\emph{leasing}), limita la producción, garantiza la devolución del dinero si el precio baja o hace el bien menos duradero: obsolescencia, nuevas ediciones que acaban con el mercado de segunda mano, como en los libros de texto \parencite{Bulow1982Durable}.
\end{itemize}

\paragraph{c) Evidencia empírica}

\begin{itemize}
\item El precio de lanzamiento de los teléfonos y las consolas, que baja a los pocos meses, y las primeras ediciones en tapa dura de los libros.
\end{itemize}

\subsubsection{Eficiencia y bienestar de la discriminación}\label{sec:eficiencia-bienestar-discriminacion}

\begin{itemize}
\item El juicio depende del grado: el primero es eficiente pero deja a los consumidores sin excedente; el segundo distorsiona la cantidad del tipo bajo; el tercero solo mejora el excedente total si aumenta la cantidad. La lección general: cuanta más información e instrumentos tiene el monopolista, más excedente extrae, y eso no siempre empeora la eficiencia.
\item Una vez levantado el supuesto de precio único, pasamos al de unos costes sin ventajas de tamaño: el monopolio que nace de la tecnología.
\end{itemize}

% --------------------------------------------------------------------------
\subsection{Monopolio natural}\label{sec:monopolio-natural}

\begin{itemize}
\item Una vez vistos los monopolios que nacen de la información, pasamos al que nace de la tecnología. La \emph{conclusión anticipada}: aquí el monopolio es la forma más barata de producir, y el problema deja de ser cómo evitarlo para ser cómo regularlo.
\end{itemize}

\subsubsection{Monopolio natural uniproducto}\label{sec:monopolio-natural-uniproducto}

\paragraph{Existencia: subaditividad de costes}

\begin{itemize}
\item Un mercado es un \concepto{monopolio natural} si la función de costes es \concepto{subaditiva} en la producción relevante, es decir, si para todo reparto de $q$ entre varias empresas
  \begin{equation}
  C\Big(\sum_{j=1}^{n}q_j\Big)<\sum_{j=1}^{n}C(q_j) . \label{eq:subaditividad}
  \end{equation}
  Es más barato que toda la producción la haga una sola empresa: la eficiencia productiva exige un único productor. Es la definición de monopolio natural, no un supuesto adicional \parencite{Baumol1977Proper} \vertema{3.A.12}.
\item Con un producto, las \emph{economías de escala} en todo el tramo son condición \emph{suficiente}, pero no necesaria: hay subaditividad también con el coste medio ya creciente, entre el mínimo del CMe ($q^0$) y el punto en que 2 empresas empiezan a ser más baratas que una ($q^1$). En el gráfico~\ref{fig:subaditividad-sostenibilidad}, $\mathrm{CMe}_1$ es el coste medio de una empresa y $\mathrm{CMe}_{1+2}$ el de 2 empresas que se reparten la producción a partes iguales.
\item Es un concepto relativo a la demanda: el monopolio natural es una configuración de costes \emph{y} de demanda.\footnote{Un monopolio natural tradicional surge cuando la demanda es relativamente pequeña y la tecnología exige grandes costes fijos ligados a infraestructuras físicas (el cableado de la telefonía, las tuberías del agua). Con el tiempo puede dejar de ser sostenible, o incluso de existir, por la expansión de la demanda o por un cambio técnico que haga innecesaria parte de esas inversiones.}
\item Ejemplos: las redes (transporte y distribución de electricidad y de gas, abastecimiento de agua, vías férreas), no los servicios que circulan por ellas, que pueden prestarse en competencia.
\end{itemize}

\begin{figure}[H]
\centering
\caption{Subaditividad y sostenibilidad del monopolio natural uniproducto}\label{fig:subaditividad-sostenibilidad}
\grafico[0.62\textwidth]{subaditividad-sostenibilidad}
\fuente{elaboración propia a partir de \textcite{BaumolPanzarWillig1982Contestable}.}
\end{figure}

\paragraph{Sostenibilidad}

\begin{itemize}
\item Que sea más barato producir con una empresa no garantiza que la empresa pueda mantenerse sola. Un monopolio natural es \concepto{sostenible} si existe un precio que le permite cubrir costes atendiendo toda la demanda y al que ningún entrante, con la misma tecnología, puede obtener beneficios vendiendo a un precio igual o menor una parte de ella \parencites{PanzarWillig1977Free}{BaumolBaileyWillig1977Weak}.
\item Depende de dónde corte la demanda al coste medio (gráfico~\ref{fig:subaditividad-sostenibilidad}):
  \begin{itemize}
  \item si lo corta en su tramo decreciente, antes de $q^0$, es \emph{sostenible}: para atraer clientes, el entrante tendría que cobrar menos que el coste medio que tendría con una cantidad menor, y perdería dinero;
  \item si lo corta entre $q^0$ y $q^1$, es un monopolio natural pero, si el incumbente no puede reaccionar y debe atender toda la demanda a su precio, \emph{no es sostenible}: un entrante puede producir $q^0$, en el mínimo del coste medio, y vender algo más barato que el incumbente con beneficio;
  \item si el incumbente puede reaccionar, la sostenibilidad depende de si la demanda queda por debajo del coste medio que tendrían 2 empresas repartiéndose el mercado, $\mathrm{CMe}_{1+2}$: si queda por debajo, el entrante anticipa que, tras su entrada, las 2 tendrían pérdidas, y no entra \parencite{Segura1986Analisis}.\footnote{Con la demanda $D^{*}$, que corta al $\mathrm{CMe}_1$ entre $q^0$ y $q^1$ y queda siempre por debajo de $\mathrm{CMe}_{1+2}$, el incumbente produce $q^{*}$ a $p^{*}=\mathrm{CMe}_1(q^{*})$, con beneficio nulo. Si entra otra empresa, repartirse el mercado no es rentable, porque el precio quedaría por debajo de $\mathrm{CMe}_{1+2}$, y al incumbente solo le queda bajar el precio (y la cantidad) hasta el mínimo del coste medio, $p^0$ y $q^0$, donde el entrante tendría pérdidas: el entrante lo anticipa y no entra. Con una demanda $D'$ que en algún tramo supera a $\mathrm{CMe}_{1+2}$, el incumbente empieza bajando el precio, pero a partir de cierto punto le conviene acomodar la entrada y repartirse el mercado a un precio por encima de $\mathrm{CMe}_{1+2}$, con beneficios, en lugar de seguir luchando con beneficio nulo: el entrante lo anticipa y entra, y las 2 empresas producen donde el beneficio conjunto es máximo. La mayoría de los manuales suponen, como \autor{Baumol}, \autor{Panzar} y \autor{Willig}, que el incumbente no reacciona y mantiene su precio, y por eso concluyen que el monopolio natural no es sostenible en todo el tramo entre $q^0$ y $q^1$. La secuencia completa está en el anexo~\ref{anexo:sostenibilidad}.}
  \end{itemize}
\item \emph{Relación con los mercados impugnables}: si la entrada y la salida son libres y sin costes hundidos, la sola amenaza de entrada obliga al monopolio natural sostenible a cobrar el coste medio, sin beneficios extraordinarios \parencite{BaumolPanzarWillig1982Contestable} \vertema{3.A.15}.
\end{itemize}

\subsubsection{Monopolio natural multiproducto}\label{sec:monopolio-natural-multiproducto}

\begin{itemize}
\item Con varios productos, la definición sigue siendo la \concepto{subaditividad} de $C(\vect{w},\vect{q})$: ningún reparto de $\vect{q}$ entre varias empresas es más barato que producirlo todo en una. Las \concepto{economías de alcance},
  \begin{equation}
  C(q_1,q_2)<C(q_1,0)+C(0,q_2),
  \end{equation}
  es decir, producir los 2 bienes en una empresa es más barato que en 2 empresas especializadas, son condición \emph{necesaria}, no suficiente; una condición suficiente es que, además, el coste incremental medio de cada producto sea decreciente \parencite{BaumolPanzarWillig1982Contestable} \vertema{3.A.12}.
\item \emph{Sostenibilidad}. Exige, además, precios \concepto{libres de subsidios cruzados} \parencite{Faulhaber1975Cross}: el monopolista no puede vender unos bienes caros para financiar otros por debajo de su coste. Un vector de precios está libre de subsidios cruzados si pasa 2 pruebas:
  \begin{itemize}
  \item \emph{prueba del coste aislado}: ningún subconjunto de productos genera ingresos mayores que lo que costaría producirlo por separado; si no, un entrante produciría solo esos bienes y los vendería más baratos;
  \item \emph{prueba del coste incremental}: ningún subconjunto genera ingresos menores que lo que cuesta añadirlo a la producción del resto; si no, los demás productos lo están financiando.
  \end{itemize}
  Es una condición \emph{necesaria} de la sostenibilidad, no suficiente; si la empresa cubre exactamente sus costes, las 2 pruebas son equivalentes.
\end{itemize}

\subsubsection{Regulación del monopolio natural}\label{sec:regulacion-monopolio-natural}

\paragraph{a) Idea}

\begin{itemize}
\item Si el monopolio natural es sostenible, no conviene desmantelarlo, porque se perdería la eficiencia productiva, pero tampoco dejarlo libre, porque fijaría el precio de monopolio. La regulación busca el mejor equilibrio entre la \emph{eficiencia asignativa} y la \emph{viabilidad financiera} de la empresa.
\end{itemize}

\paragraph{b) Modelo}

\subparagraph{i) Supuestos}

\begin{itemize}
\item Un regulador benevolente, que maximiza el excedente total, e informado: conoce los costes y la demanda. La regulación con información asimétrica y la teoría positiva de la regulación están en el anexo~\ref{anexo:regulacion} \vertema{3.A.20}.
\end{itemize}

\subparagraph{ii) Desarrollo}

\begin{itemize}
\item \concepto{Primer óptimo} (\emph{first best}): \emph{precio igual al coste marginal con subvención}.\footnote{Hablamos de primer y segundo óptimo, y no de óptimos «de primer y segundo orden», para no confundirlos con las condiciones de primer y segundo orden de los problemas de optimización \vertema{3.A.22}.}
  \begin{itemize}
  \item Con $p=\mathrm{CMg}$ la asignación es eficiente \parencite{Hotelling1938General}. Pero, si la demanda corta en el tramo decreciente del coste medio, $\mathrm{CMg}<\mathrm{CMe}$ y la empresa tiene pérdidas iguales a $(\mathrm{CMe}-\mathrm{CMg})\,q$ (gráfico~\ref{fig:regulacion-monopolio-natural}), que hay que cubrir con una subvención de suma fija. Es lo mismo que impide la competencia perfecta con rendimientos crecientes: una empresa precio-aceptante pierde dinero \vertema{3.A.21}.
  \item Problemas: los impuestos que financian la subvención distorsionan, y cada euro de subvención cuesta $1+\lambda$ euros de bienestar, con $\lambda>0$ el \emph{coste sombra de los fondos públicos}; la subvención garantizada quita incentivos a reducir costes, y redistribuye de los contribuyentes a los usuarios.\footnote{Hay, por tanto, un \emph{trade-off} entre la distorsión de los impuestos y la de los precios; un problema de eficiencia dinámica, porque la empresa sabe que sus pérdidas se cubrirán, y uno redistributivo, entre quienes pagan los impuestos y quienes usan el servicio. Además, en la práctica el regulador rara vez conoce bien los costes de la empresa.}
  \item Si la demanda corta en el tramo creciente del coste medio, $p=\mathrm{CMg}$ deja beneficios, que se pueden gravar con un impuesto de suma fija (anexo~\ref{anexo:regulacion}).
  \end{itemize}
\item \concepto{Segundos óptimos} (\emph{second best}): cubrir los costes sin subvención.
  \begin{itemize}
  \item \emph{Precio igual al coste medio}: beneficio nulo y la mayor cantidad compatible con no tener pérdidas, pero con pérdida de eficiencia (se producen menos unidades de las eficientes), y el regulador necesita conocer los costes, que la empresa tiende a inflar; además, si el precio se revisa para cubrir costes, la empresa deja de tener incentivos a reducirlos.
  \item \emph{Tarifa en dos partes} \parencite{Coase1946Marginal}: una cuota fija $A$ y un precio por unidad $p=\mathrm{CMg}$, como en la discriminación de segundo grado (\ref{sec:discriminacion-segundo-grado}). Con $n$ usuarios y coste fijo $\mathrm{CF}$, la cuota $A=\mathrm{CF}/n$ cubre las pérdidas sin distorsionar el margen.\footnote{Si la cuota no expulsa a nadie, funciona como una transferencia de suma fija del consumidor a la empresa: no cambia ninguna decisión en el margen y la asignación es eficiente. El coste de la tarifa en dos partes es distributivo, salvo que la cuota expulse a consumidores que valoran el bien por encima de su coste marginal.} Funciona si la cuota no expulsa a nadie; si expulsa a usuarios que valoran el servicio por encima de su coste, conviene bajar la cuota y subir el precio por encima del coste marginal \parencite{Oi1971Disneyland}, o cobrar cuotas distintas según la capacidad de pago (por ejemplo, con bonos sociales), que es discriminar en la parte fija.\footnote{La cuota no tiene por qué ser igual para todos: para no expulsar a los usuarios de menos renta se les puede cobrar menos, financiándolo con cuotas mayores al resto; y, por eficiencia, conviene cargar más la cuota a quienes tienen una demanda de acceso menos elástica, la lógica de \autor{Ramsey} aplicada a la parte fija \parencite{ViscusiHarringtonVernon2005Economics}. Las 2 cosas pueden chocar: los usuarios de demanda de acceso más inelástica suelen ser los de menor renta, porque los servicios de red son de primera necesidad.}
  \item \emph{Tarifas de punta y valle}: si la demanda varía en el tiempo (por horas o estaciones), el producto no se puede almacenar y la capacidad se fija para la punta, la tarifa eficiente cobra en el valle solo el coste variable, $p_v=b$, y en la punta también el coste de la capacidad, $p_p=b+\beta$ \parencites{Boiteux1949Tarification}{Steiner1957Peak}. El desarrollo y la aplicación a la factura eléctrica están en el anexo~\ref{anexo:regulacion}.
  \item \emph{Precios de \autor{Ramsey}-\autor{Boiteux}}, para el monopolio natural multiproducto: el regulador maximiza el excedente total sujeto a que la empresa no tenga pérdidas,
    \begin{equation}
    \optim{\max}{\vect{p}}{\mathrm{ET}(\vect{p})=\mathrm{EC}(\vect{p})+B(\vect{p})}{B(\vect{p})\geq 0}
    \end{equation}
    y, con demandas independientes, obtiene \parencites{Ramsey1927Contribution}{Boiteux1956Gestion} (anexo~\ref{anexo:ramsey}):
    \begin{equation}
    \frac{p_i-\mathrm{CMg}_i}{p_i}=\frac{\lambda}{1+\lambda}\cdot\frac{1}{|\varepsilon_{q_i,p_i}|}\qquad\forall i, \label{eq:ramsey-boiteux}
    \end{equation}
    con $\lambda$ el multiplicador de la restricción presupuestaria: el bienestar que se ganaría si se exigiera un euro menos de beneficio.
  \end{itemize}
\end{itemize}

\begin{figure}[H]
\centering
\caption{Regulación del monopolio natural: pérdidas con $p=\mathrm{CMg}$ y beneficio nulo con $p=\mathrm{CMe}$}\label{fig:regulacion-monopolio-natural}
\grafico[0.6\textwidth]{regulacion-monopolio-natural}
\fuente{elaboración propia a partir de \textcite{Braeutigam1989Optimal}.}
\end{figure}

\subparagraph{iii) Implicaciones}

\begin{itemize}
\item La regla de \autor{Ramsey}-\autor{Boiteux} es la \emph{regla de \autor{Lerner} amortiguada}. El número de \autor{Ramsey}, $\lambda/(1+\lambda)\in[0,1)$, \emph{crece} con $\lambda$: si la restricción no aprieta ($\lambda=0$), $p_i=\mathrm{CMg}_i$; cuanto más aprieta (más coste fijo que cubrir), mayores son los márgenes, y si apretara sin límite ($\lambda\to\infty$), cada precio sería el de monopolio.
\item Entre ambos extremos, el regulador recarga más los bienes de demanda menos elástica, como el monopolista que discrimina en tercer grado (\ref{sec:discriminacion-tercer-grado}), porque ahí subir el precio distorsiona menos la cantidad, pero solo lo justo para cubrir costes. Es la misma regla de la elasticidad inversa de la imposición óptima \vertema{4.B.11}.
\item Con transferencias públicas que cuestan $1+\lambda$, la regla es la misma, con $\lambda$ el coste sombra de los fondos públicos \parencite[cap.~3]{LaffontTirole1993Theory}: por eso se usa la misma letra en los 2 casos. Con demandas dependientes aparecen las elasticidades cruzadas, como en el monopolio multiproducto (\ref{sec:monopolio-multiproducto}).
\end{itemize}

\paragraph{c) Valoración}

\begin{itemize}
\item La regla de \autor{Ramsey}-\autor{Boiteux} exige conocer los costes y las elasticidades, y puede ser regresiva: los bienes de demanda inelástica suelen ser de primera necesidad y pesan más en la cesta de los hogares de menos renta.
\item Todos estos esquemas suponen que el regulador conoce los costes y la demanda y que solo le mueve el bienestar. Cuando no es así (información asimétrica, captura del regulador), aparecen la regulación de la tasa de rentabilidad, los precios máximos, los menús de contratos y la competencia por comparación (anexo~\ref{anexo:regulacion}) \vertema{3.A.20}.
\item \emph{Eficiencia y bienestar}: el monopolio natural es eficiente en los costes y, bien regulado, puede acercarse a la eficiencia asignativa; el coste de lograrlo depende de la información del regulador. Una vez levantado el supuesto de costes sin ventajas de tamaño, pasamos al de un solo producto y una sola planta.
\end{itemize}

% --------------------------------------------------------------------------
\subsection{Producción conjunta}\label{sec:produccion-conjunta}

\begin{itemize}
\item Una vez levantados los supuestos de precio único y de costes sin ventajas de tamaño, levantamos el de un solo bien y una sola planta. La \emph{conclusión anticipada}: el monopolista que produce varios bienes \emph{internaliza} las relaciones entre ellos, en la demanda y en los costes, y eso cambia sus precios respecto a los de varios monopolistas independientes.
\end{itemize}

\subsubsection{Monopolio multiproducto}\label{sec:monopolio-multiproducto}

\begin{itemize}
\item Con $\vect{q}=(q_1,\dots,q_n)$, demandas $q_i(\vect{p})$ y costes $C(\vect{q})$, el resultado depende de si las demandas y los costes de los bienes están relacionados. Hay 4 casos:
  \begin{enumerate}[label=\roman*)]
  \item \emph{Demandas independientes y costes separables}: cada bien es un monopolio independiente, con $\mathcal{L}_i=1/|\varepsilon_{q_i,p_i}|$, como en el monopolio uniproducto.
  \item \emph{Demandas dependientes y costes separables}: el margen de cada bien tiene en cuenta lo que su precio hace vender de los demás (anexo~\ref{anexo:multiproducto}):
    \begin{equation}
    \mathcal{L}_i=\frac{1}{|\varepsilon_{q_i,p_i}|}+\sum_{j\neq i}\frac{(p_j-\mathrm{CMg}_j)\,q_j}{p_i\,q_i}\cdot\frac{\varepsilon_{q_j,p_i}}{|\varepsilon_{q_i,p_i}|} . \label{eq:lerner-multi}
    \end{equation}
    \begin{itemize}
    \item \emph{Sustitutivos} ($\varepsilon_{q_j,p_i}>0$): subir $p_i$ desvía ventas hacia $j$, que también es suyo, así que fija precios \emph{más altos} que varios monopolistas independientes: los consumidores no pueden «escapar».\footnote{Podría explicar, por ejemplo, por qué un supermercado con tiendas físicas que pasa a vender también en internet fijaría precios más altos en línea que un supermercado que solo vende en internet.}
    \item \emph{Complementarios} ($\varepsilon_{q_j,p_i}<0$): bajar $p_i$ hace vender más de $j$, así que fija precios \emph{más bajos}, y su índice de \autor{Lerner} es menor que la inversa de la elasticidad propia; puede incluso vender uno por debajo de su coste marginal, con un índice de \autor{Lerner} negativo (consolas e impresoras baratas que se rentabilizan con los juegos y los cartuchos). Es la misma lógica que la de la doble marginalización (\ref{sec:monopolio-bilateral}).
    \end{itemize}
  \item \emph{Demandas independientes y costes no separables}: cuando producir un bien abarata el otro (\concepto{complementariedad de costes}, $\partial^2C/\partial q_1\partial q_2<0$, que implica economías de alcance) \vertema{3.A.12}, las condiciones de primer orden, $\mathrm{IMg}_i(q_i)=\partial C(\vect{q})/\partial q_i$, llevan a producir más de los 2 bienes que 2 empresas separadas, porque cada unidad de un bien reduce el coste marginal del otro.\footnote{Es distinto del \emph{aprendizaje}, en el que la producción de hoy abarata la de mañana: $q_1$ y $q_2$ son el mismo bien en 2 periodos. Con un factor de descuento $\beta$ y un coste de mañana $c_2(q_1)$, con $c_2'<0$, la condición de primer orden del primer periodo es $\mathrm{IMg}_1(q_1)=c_1+\beta\,c_2'(q_1)\,q_2<c_1$: el monopolista produce hoy más que uno miope, y puede vender incluso por debajo del coste marginal corriente, porque cada unidad le enseña a producir más barato. A cambio, traslada su poder de mercado al futuro. El ejemplo clásico es el de la industria aeronáutica de \textcite{Wright1936Factors} \vertema{3.A.12}.}
  \item \emph{Demandas dependientes y costes no separables}: se combinan los 2 efectos anteriores; los márgenes de \eqref{eq:lerner-multi} se calculan con los costes marginales $\partial C(\vect{q})/\partial q_i$, que dependen de la producción de los demás bienes.
  \end{enumerate}
\item En todos los casos, la condición de que el ingreso marginal de cada bien iguale su coste marginal se cumple; lo que cambia es que ambos dependen de lo que se produce de los demás. Un precio por debajo del coste marginal \emph{corriente} solo aparece con bienes complementarios en la demanda o con aprendizaje; no hay que llamarlo \emph{subsidio cruzado}, un término que se reserva para los precios que no pasan las pruebas de \textcite{Faulhaber1975Cross} (\ref{sec:monopolio-natural-multiproducto}).
\end{itemize}

\subsubsection{Producción conjunta en proporciones fijas}\label{sec:produccion-conjunta-proporciones}

\begin{itemize}
\item Hay bienes que salen \emph{juntos} de un mismo proceso y en una proporción fija, como la carne y la piel de una res: es la \concepto{oferta conjunta} de \textcite[libro~V, cap.~6]{Marshall1890Principles}. No hay un coste propio de cada bien, sino del proceso.
\item El monopolista produce hasta que lo que le aportan los 2 bienes juntos iguala el coste marginal del proceso: $\mathrm{IMg}_1(q)+\mathrm{IMg}_2(q)=\mathrm{CMg}(q)$.
\item Si a esa cantidad el ingreso marginal de uno de ellos es negativo, le conviene \emph{desechar} parte de ese bien en lugar de venderlo todo. El modelo completo, con su gráfico, está en el anexo~\ref{anexo:proporciones-fijas}.
\end{itemize}

\subsubsection{Monopolio multiplanta}\label{sec:monopolio-multiplanta}

\begin{itemize}
\item Si el monopolista tiene 2 plantas con costes $C_1(q_1)$ y $C_2(q_2)$, resuelve
  \begin{equation}
  \optim{\max}{q_1,q_2}{p(q_1+q_2)\,(q_1+q_2)-C_1(q_1)-C_2(q_2)}{q_1\geq 0\\ q_2\geq 0}
  \end{equation}
  y, en las plantas que producen, las condiciones de primer orden son $\mathrm{IMg}(Q)=\mathrm{CMg}_1(q_1)=\mathrm{CMg}_2(q_2)$, con $Q=q_1+q_2$.
\item Son 2 decisiones: \emph{cuánto} producir, igualando el IMg con el CMg de la empresa, que es la suma horizontal de los CMg de sus plantas, y \emph{cómo repartirlo}, igualando los costes marginales de las plantas, que es minimizar el coste \vertema{3.A.12}. Con costes marginales constantes y capacidades limitadas, llena primero la planta más barata, como en el orden de mérito del mercado eléctrico. Es el mismo problema que el de un cártel que maximiza el beneficio conjunto \vertema{3.A.19}.
\end{itemize}

\subsubsection{Eficiencia y bienestar de la producción conjunta}\label{sec:eficiencia-bienestar-produccion}

\begin{itemize}
\item Frente a varios monopolistas independientes, el monopolio multiproducto es mejor para los consumidores si los bienes son complementarios (baja los precios) y peor si son sustitutivos (los sube); el ahorro de costes de la producción conjunta y del reparto entre plantas es una ganancia de eficiencia productiva.
\item \emph{En suma}, discriminar, producir con costes subaditivos o vender varios productos son variaciones de la misma regla de la elasticidad inversa. En todos estos casos el poder de mercado estaba del lado del vendedor; pasamos a los mercados en que está del lado del comprador o en los 2 lados.
\end{itemize}

% ==========================================================================
\section{Otras estructuras de mercado\newline relacionadas con el monopolio}\label{sec:otras-estructuras-mercado}
\minutos{7}

\begin{notaopositor}
En el cante: el monopsonio con su problema, el gasto marginal y la segunda función clave (el espejo del índice de \autor{Lerner}), su gráfico con el equilibrio y la pérdida de bienestar, por qué no hay paro involuntario, las isocostes curvas en una frase y el salario mínimo por tramos sobre el mismo gráfico, con la evidencia de \autor{Card} y \autor{Krueger} y el SMI (\tiempo{3,5}); el monopolio bilateral, por capas sobre el mismo gráfico: la solución del monopsonista, la del sindicato que actúa como monopolista, el intervalo de indeterminación de \autor{Bowley}, el producto de \autor{Nash} escrito y su resultado (cantidad eficiente, reparto según el poder de negociación), \autor{Rubinstein} en una frase y la doble marginalización en otra (\tiempo{3,5}). La tabla de los 4 casos, el monopsonio dinámico, el cálculo de la doble marginalización y la derivación de la solución de \autor{Nash} (anexo~\ref{anexo:nash}) son de profundización.
\end{notaopositor}

\begin{itemize}
\item Hasta aquí el poder de mercado estaba del lado del vendedor. Por último, cambia de lado y, después, está en los 2 lados a la vez. La \emph{conclusión anticipada} del bloque: el monopsonio es el espejo del monopolio (compra menos y paga menos de lo que vale el factor) y, cuando las 2 partes tienen poder, el precio queda indeterminado y lo decide la negociación.
\end{itemize}

% --------------------------------------------------------------------------
\subsection{Monopsonio}\label{sec:monopsonio}

\paragraph{a) Idea}

\begin{itemize}
\item Hay \concepto{monopsonio} cuando hay un único comprador o, en la práctica, uno con poder para influir en el precio de lo que compra. El término lo introduce Joan \textcite{Robinson1933Economics}. Es frecuente en los mercados de factores: la única gran empresa de una comarca, un hospital público frente al personal de enfermería, una gran cadena de distribución frente a pequeños productores.
\item Así como el monopolista sabe que, para vender más, debe bajar el precio de todas las unidades que vende, el monopsonista sabe que, para comprar más, debe \emph{subir} el precio de todas las unidades que compra.
\end{itemize}

\paragraph{b) Modelo}

\subparagraph{i) Supuestos}

\begin{itemize}
\item Una empresa precio-aceptante en el mercado del bien, a un precio $p$ (supuesto que se levanta abajo), con producción $q=F(L)$ y el resto de factores fijos.
\item Es la única demandante de trabajo en su mercado y se enfrenta a una \concepto{oferta de trabajo} creciente: para contratar $L$ trabajadores tiene que pagar a todos el salario $w_L(L)$, con $w_L'(L)>0$ y elasticidad $\varepsilon_{L,w_L}=\frac{dL}{dw_L}\frac{w_L}{L}>0$. Paga el mismo salario a todos, sin discriminación salarial.
\end{itemize}

\subparagraph{ii) Desarrollo}

\begin{itemize}
\item \emph{Problema del monopsonista}:
  \begin{equation}
  \optim{\max}{L}{B(L)=p\,F(L)-w_L(L)\,L}{L\geq 0} \label{eq:problema-monopsonio}
  \end{equation}
\item \emph{Condición de primer orden}: contrata hasta que el \concepto{valor del producto marginal} del trabajo, $\mathrm{VPMg}_L=p\cdot\mathrm{PMg}_L$, iguala su \concepto{gasto marginal}, $\mathrm{GMg}_L=w_L(L)+w_L'(L)\,L$: contratar un trabajador más cuesta su salario más la subida que hay que dar a todos los demás. Escrito con la elasticidad de la oferta, y reordenando, se obtiene el \concepto{índice de poder de monopsonio}, el espejo del índice de \autor{Lerner}:
  \funcionclave[eq:monopsonio]{\mathrm{VPMg}_L=\mathrm{GMg}_L=w_L(L)\Big(1+\frac{1}{\varepsilon_{L,w_L}}\Big)\ \Longleftrightarrow\ \frac{\mathrm{VPMg}_L-w_L}{w_L}=\frac{1}{\varepsilon_{L,w_L}}}
  La brecha entre lo que aporta el último trabajador y lo que cobra es mayor cuanto \emph{menos elástica} es la oferta de trabajo que afronta la empresa. Con oferta perfectamente elástica, $\varepsilon_{L,w_L}\to\infty$, la brecha desaparece y volvemos a la empresa precio-aceptante, $\mathrm{VPMg}_L=w_L$.
\item \emph{Interpretación gráfica} (gráfico~\ref{fig:monopsonio}): la curva de gasto marginal queda \emph{por encima} de la oferta de trabajo (con oferta lineal, con la misma ordenada en el origen y doble pendiente). La empresa elige $L^M$ donde el $\mathrm{GMg}_L$ corta al $\mathrm{VPMg}_L$, y paga el salario que la \emph{oferta} da para esa cantidad, $w^M=w_L(L^M)$. El equilibrio competitivo, $L^c$ y $w^c$, está donde el $\mathrm{VPMg}_L$ corta a la oferta.
\item \emph{No hay curva de demanda de trabajo}, por la misma razón que el monopolista no tiene curva de oferta: el empleo depende de la forma de la oferta de trabajo, no solo de su altura.
\item \emph{En los costes}: si el precio de un factor depende de cuánto se compra, las isocostes $\overline{C}=w_L(L)\,L+w_K\,K$ dejan de ser rectas. Su pendiente es $-\mathrm{GMg}_L/w_K$, cada vez mayor en valor absoluto, y la combinación de mínimo coste iguala la relación marginal de sustitución técnica al cociente de \emph{gastos marginales}, $\mathrm{RMST}_{LK}=\mathrm{GMg}_L/w_K$, no al de los precios \vertema{3.A.12}.
\item \emph{Si la empresa es además monopolista en el mercado del bien}, a la izquierda queda el \concepto{ingreso del producto marginal}, $\mathrm{IPMg}_L=\mathrm{IMg}\cdot\mathrm{PMg}_L$, menor que el VPMg: $\mathrm{IPMg}_L=w_L(L)\big(1+1/\varepsilon_{L,w_L}\big)$. Se suman las 2 brechas, la del poder en el bien y la del poder en el factor (tabla~\ref{tab:cuatro-casos}).
\end{itemize}

\begin{figure}[H]
\centering
\caption{Monopsonio en el mercado de trabajo: menos empleo y menos salario que en competencia}\label{fig:monopsonio}
\grafico[0.6\textwidth]{monopsonio}
\fuente{elaboración propia a partir de \textcite{Robinson1933Economics} y \textcite{Manning2003Monopsony}.}
\end{figure}

\begin{table}[H]
\centering
\caption{Contratación de trabajo según el poder de mercado de la empresa en el bien y en el factor}\label{tab:cuatro-casos}
\renewcommand{\arraystretch}{1.3}
\begin{tabular}{>{\centering\arraybackslash}m{0.24\textwidth}>{\centering\arraybackslash}m{0.3\textwidth}>{\centering\arraybackslash}m{0.3\textwidth}}
\toprule
\rowcolor{tceeazulpalido} & \textbf{Trabajo: competencia}\newline (oferta horizontal a $w_L$) & \textbf{Trabajo: monopsonio}\newline (oferta creciente $w_L(L)$)\\
\midrule
\textbf{Bien: competencia}\newline ($\mathrm{IMg}=p$) & \cellcolor{tceesalvia}\textcolor{tceeverde}{\textbf{1.} $\mathrm{VPMg}_L=w_L$}\newline empleo y salario máximos & \textcolor{tceeazul}{\textbf{2.} $\mathrm{VPMg}_L=\mathrm{GMg}_L$}\newline explotación monopsonística\\
\textbf{Bien: monopolio}\newline ($\mathrm{IMg}<p$) & \textcolor{tceedorado}{\textbf{3.} $\mathrm{IPMg}_L=w_L$}\newline explotación monopolística & \textcolor{tceerojo}{\textbf{4.} $\mathrm{IPMg}_L=\mathrm{GMg}_L$}\newline empleo y salario mínimos\\
\bottomrule
\end{tabular}
\fuente{elaboración propia a partir de \textcite{Robinson1933Economics}.}
\end{table}

\begin{itemize}
\item \emph{Lectura de la tabla}: con competencia en los 2 mercados (caso 1), el empleo y el salario son los más altos; con poder en los 2 (caso 4), los más bajos; con poder solo en uno (casos 2 y 3), quedan entre ambos, sin un orden claro entre ellos. En \emph{ninguno} de los 4 casos hay paro involuntario: la empresa contrata sobre la curva de oferta, y al salario pagado trabajan todos los que quieren. El paro involuntario aparece cuando el salario se fija por encima del de la oferta, por un sindicato (\ref{sec:monopolio-bilateral}) o por un salario mínimo superior al competitivo.
\end{itemize}

\subparagraph{iii) Implicaciones}

\begin{itemize}
\item \emph{Menos empleo y menos salario, no paro}: $L^M<L^c$ y $w^M<w^c$. Los trabajadores están sobre su curva de oferta: hay menos empleo, pero no hay paro involuntario.
\item \emph{Bienestar}: se pierde el triángulo entre el $\mathrm{VPMg}_L$ y la oferta de trabajo desde $L^M$ hasta $L^c$ (trabajo que vale más de lo que cuesta y no se contrata), y el rectángulo $(w^c-w^M)\,L^M$ pasa de los trabajadores a la empresa. Es la \emph{explotación monopsonística} de \textcite{Robinson1933Economics}: el trabajador cobra menos que el valor de lo que aporta en el margen.
\item \emph{Salario mínimo} (gráfico~\ref{fig:salario-minimo}). Un salario mínimo $\underline{w}$ hace horizontal la oferta de trabajo, y con ella el gasto marginal, hasta $L^S(\underline{w})$: contratar un trabajador más ya no obliga a subir el salario de los demás.
  \begin{itemize}
  \item si $\underline{w}\leq w^M$, no afecta;
  \item si $w^M<\underline{w}\leq w^c$, el empleo \emph{sube}, hasta $L^S(\underline{w})$, porque la empresa ya no gana nada contratando menos para pagar menos; el máximo de empleo se alcanza con $\underline{w}=w^c$;
  \item si $\underline{w}>w^c$, el empleo lo da el $\mathrm{VPMg}_L$, $L^D(\underline{w})<L^S(\underline{w})$, y aparece paro involuntario; aun así, el empleo es mayor que el de monopsonio mientras $\underline{w}$ no supere el valor del producto marginal en $L^M$.
  \end{itemize}
  Con monopsonio, un salario mínimo moderado puede subir a la vez los salarios y el empleo, al revés que en competencia \vertema{3.A.25}.
\end{itemize}

\begin{figure}[H]
\centering
\caption{Salario mínimo con monopsonio: el empleo sube hasta $w^c$ y baja por encima}\label{fig:salario-minimo}
\grafico[0.6\textwidth]{salario-minimo}
\fuente{elaboración propia a partir de \textcite{Manning2003Monopsony}.}
\end{figure}

\paragraph{c) Extensiones}

\begin{itemize}
\item \emph{Monopsonio dinámico} \parencite{Manning2003Monopsony}: aunque haya muchas empresas, las fricciones (los costes de buscar y de cambiar de empleo, la preferencia por un lugar) hacen que la oferta de trabajo que afronta cada una no sea perfectamente elástica: si baja el salario, no se van todos. El poder de monopsonio no exige un único empleador.
\item \emph{Monopsonio en los mercados de bienes}: los grandes distribuidores frente a los proveedores agrícolas, y el debate europeo sobre las prácticas comerciales desleales en la cadena alimentaria \vertema{3.B.39}.
\end{itemize}

\paragraph{d) Evidencia empírica}

\begin{itemize}
\item \textcite{CardKrueger1994Minimum} comparan 410 restaurantes de comida rápida de Nueva Jersey, que subió el salario mínimo de 4,25 a 5,05 dólares por hora en abril de 1992, y de Pensilvania, que no lo subió, y no encuentran indicios de que la subida redujera el empleo: un resultado coherente con el monopsonio y no con la competencia perfecta.
\item \textcite{YehMacalusoHershbein2022Monopsony} estiman que, en la planta manufacturera media de Estados Unidos, el ingreso del producto marginal del trabajo supera al salario en un 53~\%: el trabajador cobra unos 65 centavos por cada dólar que genera en el margen. \textcite{DebEtAl2022Wage} atribuyen el estancamiento de los salarios en Estados Unidos entre 1997 y 2016 sobre todo al poder de monopolio de las empresas en sus mercados de bienes (un 75~\%) y en menor medida al de monopsonio (un 25~\%).
\item En España, el salario mínimo interprofesional (SMI) de 2026 es de 1.221 euros al mes en 14 pagas (17.094 euros al año), un 3,1~\% más que en 2025 \parencite{BOE2026SMI}.
\end{itemize}

\paragraph{e) Valoración}

\begin{itemize}
\item El modelo explica por qué el salario mínimo no siempre destruye empleo y por qué el poder de mercado de las empresas puede deprimir los salarios. Su límite es que supone un único salario para todos y deja fuera la respuesta de los trabajadores: si se organizan en un sindicato, hay poder de mercado en los 2 lados. Es el monopolio bilateral.
\end{itemize}

% --------------------------------------------------------------------------
\subsection{Monopolio bilateral}\label{sec:monopolio-bilateral}

\begin{itemize}
\item Una vez estudiado el poder del comprador, pasamos al caso en que vendedor y comprador tienen poder a la vez.
\end{itemize}

\paragraph{a) Idea}

\begin{itemize}
\item Hay \concepto{monopolio bilateral} cuando un único vendedor se enfrenta a un único comprador: un sindicato que representa a todos los trabajadores frente a un empleador monopsonista, un fabricante de componentes frente al único montador que los usa, un gestor de infraestructura frente a su único usuario.
\item Cada uno querría imponer su precio y ninguno es precio-aceptante. \textcite{Bowley1928Bilateral} muestra que la teoría del monopolio y la del monopsonio dan cada una su precio y que el resultado queda \emph{indeterminado} entre ambos: no lo deciden los costes y la demanda, sino el \emph{poder de negociación} \parencite{Fellner1947Prices}. Es el mismo problema que \textcite{Edgeworth1881Mathematical} plantea con 2 individuos que intercambian: un abanico de acuerdos posibles y ningún precio de mercado que elija uno \vertema{3.A.21}.
\end{itemize}

\paragraph{b) Modelo}

\subparagraph{i) Supuestos}

\begin{itemize}
\item Los de \ref{sec:monopsonio}, con un único oferente de trabajo: un \emph{sindicato} que representa a todos los trabajadores y negocia en su nombre con el empleador monopsonista.
\item Para el sindicato, el coste de oportunidad de que trabajen $L$ personas es la desutilidad agregada del trabajo, $v(L)$, con $v'(L)=w_L(L)$: la oferta de trabajo de \ref{sec:monopsonio} mide el salario de reserva del último trabajador. El sindicato maximiza la renta de sus miembros, $U_S=w\,L-v(L)$; la empresa, su beneficio, $B=p\,F(L)-w\,L$.
\end{itemize}

\subparagraph{ii) Desarrollo}

\begin{itemize}
\item \emph{Los 2 extremos}, por capas sobre el gráfico del monopsonio (gráfico~\ref{fig:monopolio-bilateral}):
  \begin{itemize}
  \item \emph{si manda la empresa}, es la solución de monopsonio de \ref{sec:monopsonio}: $L^M$ y $w^M$;
  \item \emph{si manda el sindicato}, actúa como un monopolista que vende trabajo frente a la demanda de la empresa, $\mathrm{VPMg}_L$: fija el empleo donde su ingreso marginal, $\mathrm{IMg}_S$ (la curva marginal de la demanda de trabajo), iguala el coste de oportunidad del trabajo, la oferta $w_L(L)$, y cobra el salario que la demanda da para ese empleo, $w^S$, por encima del competitivo. Es el \emph{sindicato monopolista} de \textcite{Dunlop1944Wage}: al salario $w^S$ hay más gente dispuesta a trabajar que puestos, y hay paro involuntario.
  \end{itemize}
\item \emph{Indeterminación}: si negocian solo el salario y la empresa decide después cuánto emplea, el salario queda en algún punto del intervalo $[w^M,w^S]$, y el empleo es ineficiente salvo en $w^c$. Como con el salario mínimo, un salario negociado entre $w^M$ y $w^c$ \emph{aumenta} el empleo respecto al monopsonio: es el \emph{poder compensatorio} de \textcite{Galbraith1952American}.
\item \emph{Negociación eficiente}. Si negocian a la vez el salario y el empleo, les conviene llegar a la cantidad que maximiza el excedente conjunto, $p\,F(L)-v(L)$, y repartírselo. La solución de \textcite{Nash1950Bargaining} maximiza el producto ponderado de las ganancias de cada parte respecto de lo que obtendría sin acuerdo ($U_S^0$, $B^0$):
  \begin{equation}
  \optim{\max}{w,\,L}{\big(U_S-U_S^0\big)^{\gamma}\,\big(B-B^0\big)^{1-\gamma}}{U_S=w\,L-v(L)\\ B=p\,F(L)-w\,L} \label{eq:nash-bilateral}
  \end{equation}
  con $\gamma\in(0,1)$ el poder de negociación del sindicato. Las condiciones de primer orden (anexo~\ref{anexo:nash}) dan:
  \begin{equation}
  \mathrm{VPMg}_L(L^{*})=w_L(L^{*})\ \Rightarrow\ L^{*}=L^c,\qquad w^{*}L^{*}=v(L^{*})+U_S^0+\gamma\,\big[p\,F(L^{*})-v(L^{*})-U_S^0-B^0\big] .
  \end{equation}
  La cantidad es la eficiente y \emph{no depende} de $\gamma$; el salario reparte el excedente: cada parte recibe lo que tendría sin acuerdo más su cuota del excedente neto.
\item \emph{¿De dónde sale $\gamma$?} En el juego de ofertas alternadas de \textcite{Rubinstein1982Perfect}, las 2 partes se turnan para proponer un reparto y cada ronda sin acuerdo les cuesta, porque descuentan el futuro: se acuerda en la primera oferta y se lleva más el más paciente. Cuando el tiempo entre ofertas tiende a 0, el reparto converge a una solución de \autor{Nash} con $\gamma=\rho_B/(\rho_S+\rho_B)$, donde $\rho_S$ y $\rho_B$ son las tasas de descuento del sindicato y de la empresa \parencite{BinmoreRubinsteinWolinsky1986Nash} \vertema{3.A.14}.
\end{itemize}

\begin{figure}[H]
\centering
\caption{Monopolio bilateral sobre el gráfico del monopsonio: los extremos $w^M$ y $w^S$ y la cantidad eficiente}\label{fig:monopolio-bilateral}
\grafico[0.6\textwidth]{monopolio-bilateral}
\fuente{elaboración propia a partir de \textcite{Bowley1928Bilateral}.}
\end{figure}

\subparagraph{iii) Implicaciones}

\begin{itemize}
\item Con negociación eficiente, la cantidad es la de competencia y el poder de mercado solo decide el \emph{reparto}; si se negocia solo el precio, se pierde eficiencia. El poder en un lado puede compensar el del otro y acercar el resultado al eficiente.
\item Hay incentivos a la \emph{integración vertical}: una sola empresa elegiría la cantidad eficiente sin negociar \vertema{3.A.15}.
\end{itemize}

\paragraph{c) Extensiones}

\begin{itemize}
\item \emph{Doble marginalización} \parencite{Spengler1950Vertical}: si un monopolista vende un factor a otro monopolista que vende el bien final, y lo hace a un precio lineal, cada uno añade su margen. Integrados, venderían la cantidad de monopolio; separados, la cadena produce la mitad, cobra más y gana menos en conjunto.\footnote{Con demanda final $p=a-bQ$ y coste marginal $c$ del vendedor del factor: integrados, $Q=(a-c)/2b$ y $p=(a+c)/2$. Separados, el comprador, que paga $w$ por unidad, vende $Q=(a-w)/2b$; anticipándolo, el vendedor fija $w=(a+c)/2$, y el comprador vende $Q=(a-c)/4b$ a $p=(3a+c)/4$.} La integración, o una tarifa en dos partes entre ellos (\ref{sec:discriminacion-segundo-grado}), mejora a todos, incluidos los consumidores \vertema{3.A.15}. Es la misma lógica que la de los bienes complementarios del monopolio multiproducto (\ref{sec:monopolio-multiproducto}).
\item \emph{Negociación salarial}: si el sindicato y la empresa negocian solo el salario y la empresa decide el empleo (modelo del \emph{derecho a gestionar}), el resultado está sobre la demanda de trabajo y es ineficiente; si negocian salario y empleo, están sobre la \emph{curva de contratos}, fuera de la demanda, con más empleo: son los \emph{contratos eficientes} de \textcite{McDonaldSolow1981Wage} \vertema{3.A.27}.
\end{itemize}

\paragraph{d) Evidencia empírica}

\begin{itemize}
\item La negociación colectiva entre sindicatos y empresas, los conflictos de precios entre grandes distribuidores y fabricantes o productores agrícolas, y la negociación del precio de los medicamentos entre un comprador público dominante y un laboratorio con patente son monopolios bilaterales en los que el resultado depende de las alternativas de cada parte, es decir, de sus puntos de desacuerdo.
\end{itemize}

\paragraph{e) Valoración}

\begin{itemize}
\item La teoría clásica dejaba el monopolio bilateral indeterminado; la teoría de la negociación da una respuesta, pero desplaza la pregunta al poder de negociación y a las alternativas de cada parte, que dependen de instituciones (quién propone, qué pasa si no hay acuerdo) que el modelo de mercado toma como datos.
\item \emph{Recapitulando}: el monopsonio aplica en el mercado de factores la misma regla de la elasticidad inversa, ahora con la oferta, y con poder en los 2 lados el resultado ya no lo fijan el coste y la demanda, sino la negociación. Con ello hemos recorrido el poder de mercado del vendedor, del comprador y de los 2 a la vez, y podemos recoger las conclusiones de la exposición.
\end{itemize}

% ==========================================================================
\section*{Conclusión}
\minutos{1}

\begin{itemize}
\item \textbf{\emph{Recapitulación (ideas clave)}}:
  \begin{itemize}
  \item En conclusión, el monopolista iguala el ingreso marginal con el coste marginal y carga un margen igual a la inversa de la elasticidad de la demanda: produce menos y cobra más que la competencia, con un coste social que va del triángulo de \autor{Harberger} a los costes inflados y la búsqueda de rentas (\ref{sec:teoria-monopolio-equilibrio}).
  \item Los desarrollos adaptan esa misma regla y pueden cambiar el juicio de bienestar (\ref{sec:desarrollos-teoria-monopolio}): si discrimina, extrae más excedente (todo en el primer grado, con renta informativa y distorsión abajo en el segundo, con efectos ambiguos en el tercero); si la tecnología es subaditiva, el monopolio es la forma más barata de producir y hay que regularlo, con la regla de \autor{Ramsey}-\autor{Boiteux} como \autor{Lerner} amortiguado; y con varios bienes o plantas internaliza sus relaciones de demanda y de costes.
  \item En el mercado de factores, el monopsonio es su espejo: contrata menos y paga menos de lo que vale el trabajo, sin paro involuntario; y con poder en los 2 lados el resultado lo decide la negociación (\ref{sec:otras-estructuras-mercado}).
  \end{itemize}
\item \textbf{\emph{Relevancia}}: es la base teórica de la defensa de la competencia \vertema{3.B.39}, de la regulación de las industrias de red \vertema{3.A.20} y del debate sobre el salario mínimo cuando el empleador tiene poder de mercado \vertema{3.A.25}. El poder de mercado ha vuelto al centro del debate con la evidencia de márgenes crecientes \parencite{DeLoeckerEeckhoutUnger2020Rise} y de su relación con el estancamiento de los salarios \parencite{DebEtAl2022Wage}.
\item \textbf{\emph{Extensiones y relación con otras partes del temario}}: entre los 2 polos están la competencia monopolística \vertema{3.A.18} y el oligopolio \vertema{3.A.19}; las barreras de entrada y la amenaza de entrada, que puede disciplinar al monopolista, en la teoría de la empresa \vertema{3.A.15}; la regulación con información asimétrica y la definición del mercado relevante \vertema{3.A.20}; los fallos del sector público \vertema{3.A.23}; la negociación, en la teoría de juegos \vertema{3.A.14}, y la negociación salarial \vertema{3.A.27}.
\item \textbf{\emph{Opinión}}: el problema del monopolio no es el número de empresas, sino la ausencia de disciplina. Un monopolio impugnable o bien regulado puede ser la forma más barata de producir, y una patente, el premio de la innovación; uno protegido por barreras cuesta más que su triángulo. Por eso la política adecuada depende del origen del poder: defensa de la competencia frente a las conductas, regulación con incentivos frente a los monopolios naturales y atención al poder de compra en el mercado de trabajo.
\item \textbf{\emph{Idea final}}: quien parte y reparte se queda con la mejor parte, pero no con toda: se lo impiden la elasticidad de la demanda, lo que no sabe de sus clientes, la amenaza de entrada y, si enfrente hay otro que también reparte, la negociación. Y lo que se queda no compensa lo que la sociedad pierde.
\end{itemize}

\clearpage
\appendix

% ==========================================================================
\section{Anexo: regulación y control del monopolio}\label{anexo:regulacion}

\begin{notaopositor}
Con el título actual, la regulación y el control del poder de mercado son materia de otro tema \vertema{3.A.20}. Aquí se desarrollan los instrumentos pensados para el monopolio natural (tasa de rentabilidad, \emph{price cap}, menú de contratos, competencia por comparación), que se apoyan en \ref{sec:monopolio-natural}, para las preguntas del tribunal. En el cante basta la regulación básica de \ref{sec:regulacion-monopolio-natural} y una frase de remisión.
\end{notaopositor}

\subsection{Idea y justificación}\label{anexo:idea-justificacion}

\begin{itemize}
\item El monopolio suele generar una ineficiencia respecto a la competencia perfecta (\ref{sec:analisis-eficiencia-bienestar}), y eso puede justificar su regulación y control. La \concepto{regulación del monopolio} es la intervención de los poderes públicos para fijar el marco en que actúa, normalmente con un esquema de precios y restricciones a la entrada en el mercado.
\item Su justificación parte de que un monopolio natural puede ser, en ciertos contextos, más eficiente que un mercado con varias empresas (\ref{sec:monopolio-natural}). Pero varios factores lo alejan del óptimo: ineficiencia, menor calidad, menor producción. Y hay críticas a la regulación: puede usarse en beneficio de los propios regulados y perpetuar su posición.
\item Los modelos de la teoría de la regulación pueden clasificarse así:
  \begin{itemize}
  \item \concepto{teoría tradicional}: el regulador solo busca el bienestar social y tiene información completa, así que no hay comportamiento estratégico (\ref{anexo:teoria-tradicional-regulacion});
  \item \concepto{nueva teoría de la regulación}, que levanta esos supuestos por 2 vías:
    \begin{enumerate}[label=\alph*)]
    \item la \emph{información asimétrica} en el proceso regulador y el comportamiento estratégico, cuya obra de referencia es la de \textcite{LaffontTirole1993Theory} (\ref{anexo:nueva-teoria-regulacion});\footnote{Jean \autor{Tirole} recibió el Premio Nobel de Economía en 2014 «por su análisis del poder de mercado y la regulación» \parencite{Nobel2014Tirole}. Jean-Jacques \autor{Laffont}, su coautor, había muerto en 2004, y el Nobel no se concede a título póstumo.}
    \item la \emph{teoría positiva de la regulación} (o de los intereses privados), con los trabajos de \textcite{Stigler1971Theory}: las teorías anteriores son normativas (explican cómo debería ser la regulación), y esta busca explicar cómo se comportan de hecho quienes participan en ella (\ref{anexo:teoria-positiva-regulacion}).\footnote{George \autor{Stigler} recibió el Nobel en 1982 «por sus estudios pioneros de las estructuras industriales, el funcionamiento de los mercados y las causas y efectos de la regulación pública» \parencite{Nobel1982Stigler}.}
    \end{enumerate}
  \end{itemize}
\end{itemize}

\subsection{Teoría tradicional de la regulación}\label{anexo:teoria-tradicional-regulacion}

\subsubsection{Idea}\label{anexo:idea}

\begin{itemize}
\item Ante el poder de mercado, se puede regular por el lado de la \emph{oferta} o por el de la \emph{demanda}, fomentando la agrupación de los consumidores para mejorar el reparto del excedente. Esto último tiene grandes inconvenientes: el coste de organizar a los consumidores, sus desacuerdos y la imposibilidad de evitar que algunos viajen gratis (\emph{free-riders}). Por eso nos centramos en la regulación por el lado de la oferta.
\end{itemize}

\subsubsection{Desarrollo: un esquema por etapas}\label{anexo:desarrollo-esquema-etapas}

\begin{itemize}
\item \textcite{Braeutigam1989Optimal} propone un esquema por etapas que guía la intervención según las circunstancias de cada caso (gráfico~\ref{fig:algoritmo-regulacion}). Lo completamos con una rama para cuando el regulador no conoce los costes y la demanda.
\end{itemize}

\begin{figure}[H]
\centering
\caption{Esquema por etapas para la regulación del poder de mercado}\label{fig:algoritmo-regulacion}
\grafico[0.9\textwidth]{algoritmo-regulacion}
\fuente{elaboración propia a partir de \textcite{Braeutigam1989Optimal}.}
\end{figure}

\paragraph{Tipo de regulación según el origen del poder de mercado}

\begin{itemize}
\item \emph{Comportamientos estratégicos}: la medida adecuada es la \emph{defensa de la competencia}, que sanciona las conductas del monopolista que, en lugar de competir con mejores precios o productos, levantan barreras a la entrada, y prohíbe las concentraciones que crean monopolios.\footnote{En la Unión Europea, las normas sobre competencia necesarias para el funcionamiento del mercado interior son una de las 5 competencias exclusivas de la Unión (art.~3 del TFUE) \vertema{3.B.39}. Destacan 4 pilares. El primero, la prohibición de los acuerdos colusorios: «serán incompatibles con el mercado interior y quedarán prohibidos todos los acuerdos entre empresas, las decisiones de asociaciones de empresas y las prácticas concertadas que puedan afectar al comercio entre los Estados miembros y que tengan por objeto o efecto impedir, restringir o falsear el juego de la competencia dentro del mercado interior» (art.~101) \parencite{UE2016TFUE}. El segundo, la prohibición del abuso de posición dominante: no se penaliza la posición dominante en sí, sino su «explotación abusiva» (art.~102). El tercero, el control de las concentraciones, regulado por el Reglamento 139/2004: se declara incompatible toda concentración que obstaculice de forma significativa la competencia efectiva en el mercado interior o en una parte sustancial de él \parencite{UE2004Reglamento139}. El cuarto, el control de las ayudas de Estado (arts.~107 a 109 del TFUE).} Esta intervención debe guiarse por 2 principios:
  \begin{enumerate}
  \item no condenar de antemano las concentraciones, porque lo que parece anticompetitivo puede aumentar la eficiencia por economías de escala y de alcance o reducir los costes de transacción de los contratos incompletos \vertema{3.A.15};
  \item partir de la teoría de los mercados impugnables (contestables) \parencite{Baumol1982Contestable} \vertema{3.A.15}.
  \end{enumerate}
\item \emph{Posesión en exclusiva de recursos} (patentes, licencias): la medida adecuada es controlar los precios para evitar subidas bruscas.
\item \emph{Estrechez del mercado}: la estructura de costes y la demanda hacen óptimo concentrar la producción en una sola empresa, es decir, hay subaditividad de costes: es un monopolio natural (\ref{sec:monopolio-natural}). Lo primero es saber si es sostenible.
\end{itemize}

\paragraph{Monopolio natural no sostenible: competencia por el mercado}

\begin{itemize}
\item Si el monopolio natural \emph{no} es sostenible, la regulación puede garantizar la eficiencia productiva protegiendo su sostenibilidad (regulando la entrada para evitar una competencia destructiva) y reducir la ineficiencia asignativa.\footnote{El problema no desaparece si los poderes públicos asumen la propiedad del monopolio natural: la regulación sigue siendo necesaria, aunque sea implícita, para guiar la acción de la empresa pública.}
\item La idea de \textcite{Demsetz1968Why}: aunque solo quepa una empresa \emph{en} el mercado, puede haber competencia \emph{por} el mercado, subastando la concesión de su explotación durante un tiempo al licitante que ofrezca el precio más bajo. La subasta da competencia \emph{ex ante} y su repetición periódica, competencia \emph{ex post} \vertema{3.A.14}.\footnote{Hay una gran literatura sobre el número óptimo de participantes en la subasta, las condiciones de acceso, la selección de los oferentes y la propia técnica de la subasta \parencite{Klemperer2002WhatReallyMatters} \vertema{3.A.14}.} Problemas:
  \begin{enumerate}[label=\alph*)]
  \item la posible \emph{colusión} entre los licitantes;
  \item las \emph{ofertas poco competitivas}: los licitantes ofrecen precios basados en el coste medio más un margen, que puede ser alto porque incorpora el riesgo de que la concesión no se renueve;\footnote{La ineficiencia es mayor cuanto mayor es la diferencia entre el coste medio y el marginal, es decir, cuanto más importantes son los costes fijos.}
  \item la \emph{duración} de la concesión;\footnote{Si la concesión es \emph{corta} (menor que el plazo de amortización de los activos específicos), hay problemas sea quien sea el propietario de los activos: si invierte el sector público, la empresa no tiene incentivos a usar bien el equipo; si invierte la empresa, cargará un margen para cubrirse del riesgo de no renovar. Si es \emph{larga}, se suavizan esos problemas, pero se pierden las ventajas de la competencia potencial, y la incertidumbre sobre la demanda y los costes hace imposible diseñar y hacer cumplir un contrato completo.}
  \item el \emph{aprendizaje}: si los activos incorporan tecnologías muy especializadas, el primer concesionario tiene ventaja en la siguiente subasta.
  \end{enumerate}
  \textcite{Williamson1976Franchise} resume sus límites: con activos específicos y contratos largos e incompletos, la concesión acaba necesitando la misma supervisión que la regulación. Un ejemplo habitual es la concesión de servicios públicos locales, como la recogida de residuos.
\end{itemize}

\paragraph{Monopolio natural sostenible: eficiencia asignativa frente a viabilidad financiera}

\begin{itemize}
\item Si es sostenible, no conviene desmantelarlo (es eficiente en costes) ni dejarlo sin regular (subiría mucho el precio), y el regulador no debe exigir demasiada eficiencia asignativa: cuanto más se le obliga a comportarse como competitivo, mayores son sus pérdidas. Hay que encontrar un equilibrio entre eficiencia asignativa y viabilidad financiera.
\item \concepto{Primer óptimo: $p=\mathrm{CMg}$}. Si el precio igual al coste marginal se alcanza en el tramo \emph{decreciente} del coste medio, hay pérdidas que cubrir con una subvención de suma fija (\ref{sec:regulacion-monopolio-natural}). Si se alcanza en el tramo \emph{creciente}, $\mathrm{CMe}<\mathrm{CMg}=p$ y hay beneficios (gráfico~\ref{fig:pcmg-beneficios}), que el regulador puede extraer con un impuesto de suma fija. En este caso hay que regular a la vez el \emph{acceso} (proteger al monopolio con barreras a la entrada para evitar una competencia destructiva, con ciclos de inversión ineficientes) y el \emph{comportamiento} (extraer el beneficio extraordinario sin distorsionar, algo difícil en la práctica por la falta de información, o recurrir a los segundos óptimos).
\end{itemize}

\begin{figure}[H]
\centering
\caption{Precio igual al coste marginal en el tramo creciente del coste medio: beneficios que se pueden gravar}\label{fig:pcmg-beneficios}
\grafico[0.5\textwidth]{pcmg-beneficios}
\fuente{elaboración propia a partir de \textcite{Braeutigam1989Optimal}.}
\end{figure}

\begin{itemize}
\item \concepto{Segundos óptimos}. Ante la imposibilidad de asegurar la eficiencia asignativa sin subvención, buscan eliminar los beneficios extraordinarios:
  \begin{itemize}
  \item \emph{Criterio del coste medio}: el regulador impone $p=\mathrm{CMe}$. Limitaciones: necesita conocer los costes, y la empresa tiene incentivos a declararlos mayores (un problema de compatibilidad de incentivos); y, si el precio se revisa para cubrir costes, la empresa no tiene incentivos a reducirlos (ineficiencia dinámica).
  \item \emph{Tarifa en dos partes} \parencite{Coase1946Marginal}: cuota fija $A=\mathrm{CF}/n$ y precio por unidad $p=\mathrm{CMg}$ (\ref{sec:regulacion-monopolio-natural}). Tiene 2 problemas: supone que la demanda de acceso no depende de la cuota, cuando es previsible cierta sensibilidad, y, si la cuota es alta, puede expulsar a los consumidores de menos renta. Una solución es que la cuota sea distinta para consumidores distintos, por ejemplo según su renta: los de menos renta pagarían cuotas más bajas, compensadas con cuotas mayores al resto (\ref{sec:regulacion-monopolio-natural}).
  \item \emph{Tarifas de punta y valle} \parencites{Boiteux1949Tarification}{Steiner1957Peak}. Se aplican cuando la demanda es estacional u horaria, el producto no se puede almacenar (no hay arbitraje temporal) y la capacidad para atender la punta exige inversiones de larga maduración que dejan capacidad ociosa en el valle (turismo, energía, transporte público). Si cada unidad cuesta $b$ de operación y cada unidad de capacidad $\beta$, y la capacidad la fija la demanda de la punta (\emph{punta firme}), la tarifa eficiente es
    \begin{equation}
    p_v=b,\qquad p_p=b+\beta:
    \end{equation}
    quien usa la capacidad en la punta paga su coste, y en el valle sobra capacidad y solo se cobra el coste de operación (gráfico~\ref{fig:pico-valle}). Si con esos precios la demanda del valle superara a la de la punta (\emph{punta desplazada}), los 2 periodos comparten el coste de capacidad: con 2 periodos de igual duración, $p_p+p_v=2b+\beta$, con precios que igualan las 2 demandas a la capacidad. Así el beneficio del monopolista puede ser nulo: el coste de capacidad lo pagan quienes consumen en la punta.\footnote{Desde el 1 de junio de 2021, los consumidores domésticos de electricidad pagan los peajes y cargos con 3 periodos que fija la CNMC: punta (de 10 a 14 h y de 18 a 22 h los días laborables), llano (de 8 a 10 h, de 14 a 18 h y de 22 a 24 h) y valle (de 0 a 8 h y todas las horas de los fines de semana y los festivos nacionales), con la parte regulada del precio más cara en punta \parencite{CNMC2020Circular3}. Es una tarifa de punta y valle: quien consume en punta paga la red que se dimensiona para la punta.}
  \item \emph{Precios de \autor{Ramsey}-\autor{Boiteux}} (\ref{sec:regulacion-monopolio-natural} y anexo~\ref{anexo:ramsey}). Con demandas dependientes, además de la elasticidad propia cuentan las cruzadas, como en el monopolio multiproducto (\ref{sec:monopolio-multiproducto}): con bienes \emph{sustitutivos}, el precio regulado de cada bien es mayor, porque subirlo desvía demanda hacia el otro; con bienes \emph{complementarios}, el precio de uno puede quedar por debajo de su coste marginal para impulsar la demanda del otro, aunque los precios con subsidios cruzados no serán sostenibles (\ref{sec:monopolio-natural-multiproducto}).
  \end{itemize}
\end{itemize}

\begin{figure}[H]
\centering
\caption{Tarifas de punta y valle con punta firme}\label{fig:pico-valle}
\grafico[0.55\textwidth]{pico-valle}
\fuente{elaboración propia a partir de \textcite{Steiner1957Peak}.}
\end{figure}

\subsubsection{Valoración}\label{anexo:valoracion}

\begin{itemize}
\item La teoría tradicional supone que el regulador tiene toda la información necesaria, sin asimetrías en el proceso regulador, y que su único objetivo es el bienestar social. La nueva teoría de la regulación levanta el primer supuesto (\ref{anexo:nueva-teoria-regulacion}) y la teoría positiva, el segundo (\ref{anexo:teoria-positiva-regulacion}).
\end{itemize}

\subsection{Nueva teoría de la regulación: información asimétrica}\label{anexo:nueva-teoria-regulacion}

\begin{itemize}
\item Si la información del regulador no basta para restablecer la eficiencia asignativa o eliminar los beneficios extraordinarios, intentará al menos limitarlos. Son los \concepto{terceros óptimos}.\footnote{En la teoría de los contratos, el mejor resultado posible con las restricciones de información se llama \emph{óptimo restringido} (\emph{second best}) \vertema{3.A.13}; aquí se habla de terceros óptimos porque los segundos son los de la regulación con información completa ($p=\mathrm{CMe}$, tarifas en dos partes, \autor{Ramsey}-\autor{Boiteux}).} Todos los esquemas anteriores tienen 2 problemas: de \emph{información}, porque el regulador necesita conocer bien los costes, la demanda y las preferencias sociales, y de \emph{incentivos}, porque, si usa los costes que le declaran las empresas, estas tienen incentivos a declarar costes mayores y a no esforzarse en reducirlos \vertema{3.A.13}.
\end{itemize}

\subsubsection{Regulación de la tasa de rentabilidad y efecto \autor{Averch}-\autor{Johnson}}\label{anexo:regulacion-tasa-rentabilidad}

\begin{itemize}
\item El regulador permite a la empresa una rentabilidad máxima $s$ sobre su capital, mayor que su coste $w_K$:
  \begin{equation}
  \optim{\max}{z_K,z_L}{B=p(q)\,q-w_Kz_K-w_Lz_L}{q=F(z_K,z_L)\\ \dfrac{p(q)\,q-w_Lz_L}{z_K}\leq s}
  \end{equation}
  La restricción equivale a $B\leq(s-w_K)\,z_K$: el beneficio permitido crece con el capital.
\item \emph{Gráficamente}, en el plano $(z_K,B)$ (gráfico~\ref{fig:averch-johnson}), el beneficio sin restricción tiene forma de U invertida, porque en algún momento la productividad marginal del capital cae mientras su coste sigue constante, con un máximo en $z_K^{*}$. El beneficio permitido es la recta $(s-w_K)\,z_K$. Si $z_K^{*}$ queda por encima de la recta, no está permitido, y la empresa maximiza en $z_K'>z_K^{*}$: \emph{sobrecapitaliza}.
\item Es el \concepto{efecto \autor{Averch}-\autor{Johnson}} \parencite{AverchJohnson1962Behavior}: con la restricción activa, la empresa elige una combinación de factores con más capital que la de mínimo coste,
  \begin{equation}
  \mathrm{RMST}_{LK}=\frac{\mathrm{PMg}_L}{\mathrm{PMg}_K}=\frac{w_L}{w_K-\frac{\mu}{1-\mu}\,(s-w_K)}>\frac{w_L}{w_K},
  \end{equation}
  como si el capital fuera más barato de lo que es (anexo~\ref{anexo:aj}), lo que se traduce en costes y precios más altos.
\item \emph{Ventaja}: el beneficio y el capital invertido son datos contables que se observan y se auditan con facilidad. \emph{Inconvenientes}: la sobrecapitalización; se regula como beneficio la remuneración contable del capital, que no es el beneficio económico, y no hay incentivos a reducir costes, porque lo que se ahorra se traslada al precio en la siguiente revisión.
\item Pese a sus deficiencias, es habitual en los servicios públicos por sus escasas necesidades de información. En España, la retribución de las redes de electricidad y gas se calcula aplicando una tasa de retribución financiera a la base de activos regulada \parencite{CNMC2019Circular2}, y las tarifas de la corriente de tracción que ADIF cobra a las empresas ferroviarias se fijan a partir de sus costes, con supervisión de la CNMC \parencite{CNMC2018TarifasAdif}: una regulación por coste del servicio, emparentada con la de tasa de rentabilidad.
\end{itemize}

\begin{figure}[H]
\centering
\caption{Regulación de la tasa de rentabilidad y efecto \autor{Averch}-\autor{Johnson}}\label{fig:averch-johnson}
\grafico[0.55\textwidth]{averch-johnson}
\fuente{elaboración propia a partir de \textcite{AverchJohnson1962Behavior}.}
\end{figure}

\subsubsection{Precios máximos (\emph{price cap}, \texorpdfstring{$\mathrm{IPC}-X$}{IPC-X})}\label{anexo:precios-maximos-price}

\begin{itemize}
\item La regla nace con el informe de \textcite{Littlechild1983Regulation} para la privatización de British Telecom, y se extendió después a los demás servicios públicos (gas, agua, electricidad, aeropuertos). Aborda el dilema entre eficiencia asignativa y eficiencia productiva. El mecanismo:
  \begin{enumerate}
  \item el regulador fija un nivel inicial de ingresos que permite recuperar los «costes razonables» de la empresa;
  \item a partir de esos ingresos, y para una demanda dada, se define un precio medio;
  \item el precio (o una cesta de precios ponderada) evoluciona según la regla
    \begin{equation}
    p_{t+1}\leq p_t\,(1+\mathrm{IPC}_t-X),
    \end{equation}
    con $X$ la ganancia de productividad que se exige a la industria, revisada cada pocos años: cuanto mayor es, menos sube el precio, y las ganancias de productividad pasan a los consumidores.\footnote{«Productividad de la industria» no implica que haya varias empresas. En una autopista, por ejemplo, el monopolista puede aumentar su productividad consiguiendo que circulen más vehículos con los mismos costes o reduciendo los costes para el mismo tráfico.}
  \end{enumerate}
\item Da incentivos a reducir costes, porque el ahorro aumenta el beneficio: con un precio máximo, la única forma de ganar más es producir más barato o innovar. Para un $X$ dado, el incentivo es mayor cuanto más largo es el periodo regulatorio, pero entonces se trasladan menos ganancias a los precios: la duración óptima equilibra los 2 efectos.
\item \emph{Ventajas}: incentiva la eficiencia productiva; traslada las mejoras a los consumidores \emph{a priori}, con ajustes que anticipan mejoras previstas, y \emph{a posteriori}, al fijar el nuevo límite en cada revisión; es relativamente sencilla de administrar y facilita los reequilibrios tarifarios.
\item \emph{Inconvenientes}: hace falta mucha información para fijar $X$ y el precio inicial; al acercarse el final del periodo, el interés por reducir costes disminuye, y la empresa puede retrasar mejoras hasta el comienzo del ciclo siguiente para disfrutarlas más tiempo; y la calidad puede deteriorarse para ahorrar costes.
\item En España, la cuota de abono de Telefónica estuvo sujeta a una regla $\mathrm{IPC}-X$ hasta su liberalización en 2012 \parencite{CNMC2018CuotaAbono}. Hoy el ejemplo vivo son las tarifas de Aena: el Documento de Regulación Aeroportuaria (DORA) fija un ingreso máximo por pasajero que evoluciona con un componente X durante 5 años; el DORA 2027-2031 se aprobó en septiembre de 2026 \parencite{BOE2026DORA}. La Ley 2/2015 de desindexación, desarrollada por el Real Decreto 55/2017, desvincula del IPC las revisiones de precios del sector público \parencites{BOE2015Desindexacion}{BOE2017RD55}: por eso una regla $\mathrm{IPC}-X$ pura choca con ella y los reguladores usan índices de costes propios.
\end{itemize}

\subsubsection{Mecanismo de \autor{Loeb} y \autor{Magat}}\label{anexo:mecanismo-loeb-magat}

\begin{itemize}
\item El regulador deja al monopolista fijar el precio y le paga como subvención todo el excedente del consumidor que genera \parencite{LoebMagat1979Decentralized}. Su objetivo pasa a ser el excedente total, así que elige $p=\mathrm{CMg}$ y produce la cantidad competitiva sin que el regulador conozca sus costes.
\item \emph{Inconvenientes}: el regulador tiene que conocer la demanda para calcular la subvención, y el monopolista se queda con todo el excedente, una redistribución difícil de justificar. Se puede corregir subastando el derecho a recibir ese esquema \vertema{3.A.14}, de modo que la puja devuelva la renta al Estado.
\end{itemize}

\subsubsection{Menú de contratos de \autor{Laffont} y \autor{Tirole}}\label{anexo:menu-contratos-laffont}

\paragraph{a) Idea}

\begin{itemize}
\item El modelo de \textcite{LaffontTirole1986Using} aplica la teoría de la agencia a la regulación y la contratación pública: el principal (el regulador) contrata a un agente (la empresa) que sabe más que él sobre sus costes \parencite{LaffontTirole1993Theory}.
\end{itemize}

\paragraph{b) Modelo}

\subparagraph{i) Supuestos}

\begin{itemize}
\item El coste observado es $C=\beta-e$, con $\beta$ la eficiencia intrínseca de la empresa (información oculta: baja, $\underline{\beta}$, con probabilidad $\nu$, o alta, $\overline{\beta}$) y $e$ el esfuerzo en reducir costes (acción oculta), que le cuesta $\psi(e)$, creciente y convexa. El regulador observa $C$, pero no $\beta$ ni $e$: no sabe si un coste bajo se debe a la suerte de ser eficiente o al esfuerzo. Reembolsa el coste y paga una transferencia $t$, que cuesta a la sociedad $1+\lambda$ por euro ($\lambda$, coste sombra de los fondos públicos).\footnote{Es la misma $\lambda$ de la regla de \autor{Ramsey}-\autor{Boiteux} (\ref{sec:regulacion-monopolio-natural}): el valor social de un euro más de presupuesto. \textcite[cap.~3]{LaffontTirole1993Theory} llegan al mismo número de \autor{Ramsey}, $\lambda/(1+\lambda)$.} La empresa obtiene $U=t-\psi(e)$ y debe aceptar el contrato ($U\geq 0$).
\end{itemize}

\subparagraph{ii) Desarrollo}

\begin{itemize}
\item \emph{Información completa}: el regulador exige el esfuerzo eficiente, $\psi'(e^{*})=1$ (el último euro de esfuerzo ahorra un euro de coste), y no deja rentas.
\item \emph{Información asimétrica}: la empresa eficiente puede hacerse pasar por ineficiente esforzándose menos y obtener el mismo coste; para que no lo haga hay que dejarle una \emph{renta informativa}, que crece con el esfuerzo que se pida a la ineficiente. En el óptimo,
  \begin{equation}
  \psi'(\underline{e})=1,\qquad \psi'(\overline{e})=1-\frac{\lambda}{1+\lambda}\cdot\frac{\nu}{1-\nu}\cdot\Phi'(\overline{e})<1,
  \end{equation}
  con $\Phi(e)=\psi(e)-\psi(e-\Delta\beta)$ la renta de la empresa eficiente y $\Delta\beta=\overline{\beta}-\underline{\beta}$. No hay distorsión arriba (la eficiente hace el esfuerzo óptimo) y se pide menos esfuerzo a la ineficiente para reducir la renta de la eficiente: el mismo resultado que la discriminación de segundo grado (\ref{sec:discriminacion-segundo-grado}).
\item \emph{Puesta en práctica}: un menú de contratos lineales, $t(C)=a-b\,(C-C^a)$, con $C^a$ el coste anunciado y $b\in[0,1]$ la parte del sobrecoste que paga la empresa (la \emph{potencia} del incentivo). La empresa eficiente elige el de \emph{precio fijo} ($b=1$: se queda con todo lo que ahorre) y la ineficiente uno de baja potencia, cercano al de \emph{coste incrementado} ($b$ pequeño: se le reembolsa casi todo el coste).
\end{itemize}

\subparagraph{iii) Implicaciones}

\begin{itemize}
\item El dilema de la regulación con información asimétrica es entre \emph{incentivos} y \emph{extracción de rentas}. Los contratos de alta potencia (\emph{price cap}) dan incentivos pero dejan rentas; los de baja potencia (coste incrementado, tasa de rentabilidad), al revés.
\end{itemize}

\begin{anotaciones}
\textbf{Contratos de coste incrementado y de precio fijo} \parencite{ICEXCECO3A17}. Las medidas de regulación anteriores muestran la importancia de la información; su constatación impulsó en los años 80 la nueva teoría de la regulación, cuyas aportaciones recogen \textcite{LaffontTirole1993Theory}. La regulación se ve como un mecanismo de incentivos en el marco de la teoría de la agencia, con el regulador como principal y la empresa como agente, en situaciones de selección adversa (revelación de costes) y de riesgo moral (incentivos al esfuerzo) \vertema{3.A.13}. El regulador tiene que ponderar la eficiencia productiva, estática y dinámica, con las rentas informativas que deja a la empresa: un precio igual al coste marginal sin rentas maximiza el bienestar estático, pero no da incentivos a reducir costes, y puede convenir sacrificar parte del excedente del consumidor hoy, que pasa a ser renta informativa de la empresa, a cambio de reducciones de costes futuras.

Se puede representar con un precio lineal en el coste de la empresa,\footnote{Son soluciones a través del mercado, en las que la empresa recupera sus costes con los precios, en lugar de con transferencias, que tienen el coste sombra de los fondos públicos.}
\begin{equation}
p=b\cdot F+(1-b)\cdot C,
\end{equation}
con $F$ una cuantía fija y $C$ el coste, que depende del esfuerzo de la empresa y de factores aleatorios que el regulador no observa; $b$ es la potencia del incentivo, la misma del menú anterior. Cuanto mayor es $b$, más incentivos a reducir costes, pero esas ganancias van más al beneficio de la empresa y menos al consumidor. Los 2 casos extremos:
\begin{itemize}
\item $b=0$, \emph{coste incrementado} (\emph{cost plus}): $p=C$, el regulador cubre todos los costes. No da ningún incentivo a reducir costes, y cualquier reducción exógena de costes pasa al precio. Solo es aceptable cuando la renta que dejaría un contrato de más potencia es muy grande (empresas muy ineficientes o mucha incertidumbre sobre los costes). Es la regulación por coste del servicio que dominó los servicios públicos hasta los años 80, con precios basados en costes históricos. Un periodo regulatorio más largo amortigua la falta de incentivos: mientras el precio no cambia, la empresa gana si reduce costes.
\item $b=1$, \emph{precio fijo} (\emph{fixed price}): la empresa soporta todo el riesgo de costes y tiene los máximos incentivos a ser eficiente, pero las reducciones de costes van al beneficio, no al precio. Es la regulación por precios máximos, $\mathrm{IPC}-X$, $p_{t+1}\leq p_t\,(1+\mathrm{IPC}_t-X)$, con $X$ la mejora de productividad esperada, que traslada parte de las ganancias a los consumidores; cuanto más largo el periodo, más incentivos y menos traslado. Sus límites: hace falta mucha información para fijar $X$ y el precio inicial; si este es muy bajo, la empresa tendrá dificultades y buscará renegociar; si es muy alto, tendrá beneficios elevados, y el regulador que quiera extraerlos antes de que venza el periodo pagará un coste de credibilidad.
\end{itemize}
\end{anotaciones}

\subsubsection{Efecto trinquete y desfase regulatorio}\label{anexo:efecto-trinquete-desfase}

\begin{itemize}
\item \concepto{Efecto trinquete} (\emph{ratchet}): si el regulador usa lo que observa hoy para fijar el objetivo de mañana (por ejemplo, baja el precio permitido cuando la empresa reduce costes), la empresa sabe que cada mejora se le quitará en la siguiente revisión, y no la hace o la oculta: un coste bajo hoy se convierte en una exigencia mayor mañana \parencite{Weitzman1980Ratchet}. Sin un compromiso del regulador a largo plazo, a las empresas eficientes les conviene parecer ineficientes al principio, y el regulador no puede separar los tipos \parencites{FreixasGuesnerieTirole1985Planning}{LaffontTirole1988Dynamics}.
\item \concepto{Desfase regulatorio} (\emph{regulatory lag}): es el tiempo \emph{entre 2 revisiones} de las tarifas, durante el que el precio no cambia aunque cambien los costes. Mientras dura, la empresa se queda con lo que ahorre: da incentivos a reducir costes incluso con regulación de la tasa de rentabilidad, y cuanto más largo, más incentivos y menos traslado a los consumidores \parencites{BaumolKlevorick1970Input}{Joskow1974Inflation}.
\end{itemize}

\subsubsection{Competencia por comparación (\emph{yardstick competition})}\label{anexo:competencia-comparacion-yardstick}

\paragraph{a) Idea}

\begin{itemize}
\item Si hay varias empresas reguladas parecidas que operan en zonas distintas (distribuidoras de electricidad o de agua, hospitales), el regulador puede fijar el precio de cada una a partir de los costes de las \emph{demás}, no de los suyos \parencite{Shleifer1985Theory}. Así crea entre monopolios locales una competencia que no existe en el mercado, y resuelve el problema de la revelación de costes.
\end{itemize}

\paragraph{b) Modelo}

\subparagraph{i) Supuestos}

\begin{itemize}
\item $N$ empresas idénticas, cada una con coste unitario $c_i$ que puede reducir invirtiendo $R(c_i)$, con $R'<0$.
\end{itemize}

\subparagraph{ii) Desarrollo}

\begin{itemize}
\item El regulador fija a la empresa $i$ el precio y la transferencia
  \begin{equation}
  p_i=\bar{c}_{-i}\equiv\frac{1}{N-1}\sum_{j\neq i}c_j,\qquad T_i=\bar{R}_{-i}\equiv\frac{1}{N-1}\sum_{j\neq i}R(c_j).
  \end{equation}
  Como la empresa $i$ no influye en su propio precio, maximiza $(\bar{c}_{-i}-c_i)\,q(\bar{c}_{-i})+T_i-R(c_i)$ eligiendo $c_i$, y la condición de primer orden es $-R'(c_i)=q(\bar{c}_{-i})$: invierte hasta que el último euro ahorra un euro de coste. En el equilibrio simétrico todas eligen el coste eficiente, el precio es igual al coste y no hay rentas.
\end{itemize}

\subparagraph{iii) Implicaciones}

\begin{itemize}
\item La empresa no gana nada sesgando su declaración, porque inflar sus costes solo sube el precio de las demás, y si sus costes están por encima de la media de sus pares solo se le retribuye con arreglo a esa media, así que reducirlos le compensa: da incentivos de precio fijo sin las rentas de un \emph{price cap} mal calibrado.
\end{itemize}

\paragraph{c) Evidencia empírica}

\begin{itemize}
\item La comparación exige medir la eficiencia relativa con el análisis envolvente de datos (DEA) o con la frontera estocástica \vertema{3.A.12}. En Alemania, el regulador fija los ingresos máximos de las distribuidoras eléctricas comparando su eficiencia con los 2 métodos, se queda con el valor más favorable (con un mínimo del 60~\%) y exige eliminar la ineficiencia durante el periodo regulatorio \parencites{BNetzA2026Effizienzvergleich}{ARegV2007}; \textcite{JamasbPollitt2001Benchmarking} repasan la experiencia internacional en electricidad.
\end{itemize}

\paragraph{d) Valoración}

\begin{itemize}
\item \emph{Requisitos y problemas}: empresas comparables (si no, hay que corregir las diferencias observables, por regresión o con fronteras de eficiencia) y el riesgo de colusión entre ellas, que pueden ponerse de acuerdo en no reducir costes \parencite{BaldwinCaveLodge2012Understanding}.
\end{itemize}

\subsubsection{Evidencia empírica sobre la regulación por incentivos}\label{anexo:evidencia-empirica-regulacion}

\begin{itemize}
\item Desde los años 80, la regulación de las redes ha pasado de la tasa de rentabilidad a esquemas de incentivos (precios o ingresos máximos, comparación de eficiencia, reparto de ganancias), primero en el Reino Unido y después en la mayoría de los países europeos. \textcite{Joskow2014Incentive} concluye, en su repaso de la distribución y el transporte de electricidad, que estos esquemas han reducido los costes de operación, pero que exigen vigilar la calidad del servicio y las necesidades de inversión.
\end{itemize}

\subsection{Teoría positiva de la regulación}\label{anexo:teoria-positiva-regulacion}

\begin{itemize}
\item Las teorías anteriores son normativas: dicen cómo debería regular un regulador benevolente. La \concepto{teoría positiva} pregunta cómo se regula de hecho, y trata la regulación como un bien que se demanda y se ofrece \vertema{3.A.20}.
\item Los \concepto{grupos de presión} son conjuntos de individuos con intereses comunes que se organizan para influir en las decisiones públicas a favor de sus intereses privados (búsqueda de rentas).\footnote{La regulación en forma de política comercial (contingentes) o fiscal (exenciones) crea rentas de monopolio e incentiva a dedicar recursos a captarlas: algunas medidas regulatorias pueden camuflar intereses particulares y convertirse en una barrera de entrada injustificada.} Generan ineficiencias por 2 vías: se dedican recursos a conseguir el favor de los poderes públicos en lugar de a producir, y se toman decisiones favorables a esos grupos que no coinciden con las de la colectividad. Ejemplos: sindicatos y organizaciones empresariales.
\item \textcite{Stigler1971Theory} sostiene que, como regla general, la industria adquiere la regulación y esta se diseña y aplica sobre todo en su beneficio: es la \concepto{captura del regulador}. Las empresas reguladas son pocas, tienen mucho en juego y se organizan con facilidad; los consumidores son muchos, cada uno con poco en juego. El regulador puede no ser benevolente y responder a los intereses del grupo, lo que a su vez incentiva a las empresas a influir en el diseño de la regulación. \textcite{LaffontTirole1991Politics} formalizan la captura como un problema de agencia.
\item Es la lógica de la \concepto{acción colectiva} de \textcite{Olson1965Logic}: el beneficio de presionar es un bien público para los miembros del grupo, que tienen incentivos a viajar gratis, y el problema es menor en los grupos pequeños. Por eso solo los grupos pequeños suelen constituirse en \emph{lobbies}, y la redistribución puede ir de los grupos desorganizados a los organizados \vertema{4.B.2}.
\item \textcite{Peltzman1976Toward} añade que el regulador reparte las rentas entre productores y consumidores para maximizar sus apoyos, y \textcite{Becker1983Theory}, que la competencia entre grupos de presión favorece las intervenciones que menos eficiencia destruyen. Los recursos gastados en conseguir rentas son el coste de la búsqueda de rentas de \ref{sec:analisis-eficiencia-bienestar}.
\end{itemize}

\begin{anotaciones}
\textbf{Material complementario sobre la teoría de \autor{Stigler}}:
\begin{itemize}
\item ProMarket, «Assessing George Stigler's Economic Theory of Regulation» \parencite{ProMarket2021Stigler}: \url{https://www.promarket.org/2021/07/27/assessing-george-stiglers-economic-theory-of-regulation/}
\item Stigler Center, «Who Will Regulate the Regulators: Stigler 50 Years Later» \parencite{StiglerCenter2021Video}: \url{https://www.youtube.com/watch?v=lhD5s-4v3oM}
\end{itemize}
\end{anotaciones}

\subsection{Instituciones regulatorias}\label{anexo:instituciones-regulatorias}

\begin{itemize}
\item Los Estados combinan varios instrumentos:
  \begin{itemize}
  \item la \emph{legislación}, que establece directamente la regulación, pero choca con la falta de cualificación técnica de los legisladores, su vulnerabilidad ante los grupos organizados y el peso de las consideraciones de corto plazo;
  \item los \emph{contratos de concesión} (franquicias), que venden a las empresas el derecho a monopolizar, de modo que lo que se cobra por la venta compense las rentas del monopolista; su duración plantea un dilema: si es corta, el concesionario no invierte o lo hace pagar en el precio; si es larga, se pierde la presión de la competencia por el mercado (\ref{anexo:teoria-tradicional-regulacion});
  \item las \emph{comisiones reguladoras independientes}, con pocos miembros que aprueban las medidas regulatorias con el apoyo de equipos técnicos; más independencia de los poderes ejecutivo y legislativo suele asociarse a mejores resultados;
  \item la \emph{provisión directa} por el sector público, habitual en Europa antes de las privatizaciones de los años 80 y 90 y hoy con pocos ejemplos.
  \end{itemize}
\item En España, la Comisión Nacional de los Mercados y la Competencia (CNMC), creada por la Ley 3/2013, reúne la defensa de la competencia y la supervisión de los sectores regulados (energía, telecomunicaciones, transporte, servicios postales y audiovisuales) \parencite{Ley3_2013CNMC} \vertema{3.B.39}.
\end{itemize}

\subsection{Valoración}\label{anexo:valoracion-2}

\begin{itemize}
\item No hay regulación sin efectos indeseados: toda intervención altera los incentivos, y hay que comparar los fallos del mercado con los de la regulación, como subraya la teoría positiva \vertema{3.A.23}.
\end{itemize}

% ==========================================================================
\section{Anexo: derivaciones}\label{anexo:derivaciones}

\subsection{Condición de segundo orden del monopolista}\label{anexo:cso}

\begin{itemize}
\item De $B(q)=I(q)-C(q)$, con $I(q)=p(q)\,q$, la condición de segundo orden es
  \begin{equation}
  B''(q^M)=\frac{d^2I}{dq^2}-\frac{d^2C}{dq^2}=\mathrm{IMg}'(q^M)-\mathrm{CMg}'(q^M)<0\ \Longleftrightarrow\ \mathrm{IMg}'(q^M)<\mathrm{CMg}'(q^M).
  \end{equation}
  Las pendientes se comparan con su signo.
\item \emph{Con CMg creciente}, $\mathrm{CMg}'>0$, y con IMg decreciente, la condición se cumple siempre.
\item \emph{Con CMg decreciente}, se cumple si el CMg cae más despacio que el IMg, es decir, si corta al IMg desde abajo: a la izquierda de $q^M$ el IMg supera al CMg y a la derecha queda por debajo (gráfico~\ref{fig:cso-monopolio}, derecha). Si el CMg cae más deprisa que el IMg y lo corta desde arriba, ese punto es un mínimo del beneficio: a su derecha el IMg supera al CMg y conviene producir más.
\item A diferencia de la empresa competitiva, que solo puede maximizar en el tramo creciente del CMg (su IMg es horizontal), el monopolista puede hacerlo también en el decreciente. Si la función de beneficio es estrictamente cuasicóncava, la cantidad que maximiza el beneficio es única.\findemo
\end{itemize}

\begin{figure}[H]
\centering
\caption{Condición de segundo orden: CMg creciente (izquierda) y CMg decreciente que corta al IMg desde abajo (derecha)}\label{fig:cso-monopolio}
\grafico[0.85\textwidth]{cso-monopolio}
\fuente{elaboración propia a partir de \textcite[cap.~14]{Varian1992Microeconomic}.}
\end{figure}

\subsection{Condición de viabilidad a corto y largo plazo}\label{anexo:viabilidad}

\begin{itemize}
\item \emph{Corto plazo}. Con costes fijos $\mathrm{CF}$, $B(q^M)=\big[p(q^M)-\mathrm{CVMe}(q^M)\big]\,q^M-\mathrm{CF}$, y cerrar da $B(0)=-\mathrm{CF}$. Producir es mejor que cerrar si $B(q^M)\geq B(0)$, es decir, si $p(q^M)\geq\mathrm{CVMe}(q^M)$: las pérdidas no pueden superar el coste fijo. A corto plazo, por tanto, la condición de participación no es $B(q)\geq 0$: basta $p\geq\mathrm{CVMe}$, porque el coste fijo se paga igual.
\item \emph{Largo plazo}. Sin costes fijos, cerrar da $B(0)=0$ y producir exige $B(q^M)\geq 0$, es decir, $p(q^M)\geq\mathrm{CMe}(q^M)$: es la condición de participación $B(q)\geq 0$ del problema \eqref{eq:monopolio}.
\item \emph{Dónde se evalúa}. La condición se comprueba en $q^M$. Es equivalente a que la demanda pase por encima del coste medio en algún punto: si existe un $\hat{q}$ con $p(\hat{q})\geq\mathrm{CMe}(\hat{q})$, entonces $B(\hat{q})\geq 0$ y, como $q^M$ maximiza $B$, $B(q^M)\geq B(\hat{q})\geq 0$; y al revés, si $p(q^M)\geq\mathrm{CMe}(q^M)$, la demanda está por encima del CMe en $q^M$. Lo que no vale es exigir que el precio de monopolio supere el \emph{mínimo} del CMe, porque el monopolista no produce allí (gráfico~\ref{fig:viabilidad-monopolio}).\findemo
\end{itemize}

\begin{figure}[H]
\centering
\caption{Viabilidad del monopolista: la condición se evalúa en $q^M$, no en el mínimo del coste medio}\label{fig:viabilidad-monopolio}
\grafico[0.85\textwidth]{viabilidad-monopolio}
\fuente{elaboración propia a partir de \textcite{Koutsoyiannis1979Modern}.}
\end{figure}

\subsection{Discriminación de segundo grado con 2 tipos}\label{anexo:segundo-grado}

\begin{itemize}
\item \emph{Qué restricciones se saturan} en el problema~\eqref{eq:segundo-grado}:
  \begin{itemize}
  \item $\mathrm{PC}_2$ se deduce de $\mathrm{IC}_2$ y $\mathrm{PC}_1$: $t_2u(q_2)-T_2\geq t_2u(q_1)-T_1\geq t_1u(q_1)-T_1\geq 0$, porque $t_2>t_1$ y $u(q_1)\geq 0$. Sobra.
  \item $\mathrm{PC}_1$ se satura: si sobrara, el monopolista subiría $T_1$ y $T_2$ en la misma cuantía sin violar ninguna restricción de incentivos y ganaría más.
  \item $\mathrm{IC}_2$ se satura: si sobrara, subiría $T_2$.
  \item $\mathrm{IC}_1$ sobra en la solución: sumando las 2 restricciones de incentivos, $(t_2-t_1)\,[u(q_2)-u(q_1)]\geq 0$, es decir, $q_2\geq q_1$ (monotonía); con $\mathrm{IC}_2$ saturada, $\mathrm{IC}_1$ equivale a esa misma desigualdad, que se cumple porque $q_2>q_1$ (abajo).
  \end{itemize}
\item \emph{Pagos}: de $\mathrm{PC}_1$ y $\mathrm{IC}_2$, $T_1=t_1u(q_1)$ y $T_2=t_2u(q_2)-t_2u(q_1)+T_1=t_2u(q_2)-(t_2-t_1)\,u(q_1)$.
\item \emph{Problema reducido}:
  \begin{equation}
  \max_{q_1,q_2}\ \pi\,[t_1u(q_1)-c\,q_1]+(1-\pi)\,[t_2u(q_2)-(t_2-t_1)\,u(q_1)-c\,q_2].
  \end{equation}
\item \emph{Condiciones de primer orden}:
  \begin{equation}
  \frac{\partial}{\partial q_2}:\ (1-\pi)\,[t_2u'(q_2)-c]=0\ \Rightarrow\ t_2u'(q_2)=c;
  \end{equation}
  \begin{equation}
  \frac{\partial}{\partial q_1}:\ \pi\,[t_1u'(q_1)-c]-(1-\pi)\,(t_2-t_1)\,u'(q_1)=0\ \Rightarrow\ t_1u'(q_1)\left[1-\frac{1-\pi}{\pi}\cdot\frac{t_2-t_1}{t_1}\right]=c .
  \end{equation}
\item Si el corchete es positivo, $t_1u'(q_1)>c$: el precio marginal implícito del tipo bajo, $t_1u'(q_1)$, supera al coste marginal y $q_1$ es menor que la cantidad eficiente. Como $u'(q_1)>c/t_1>c/t_2=u'(q_2)$ y $u'$ es decreciente, $q_1<q_2$, y se cumple $\mathrm{IC}_1$. Si el corchete es nulo o negativo, $\pi\,t_1\leq(1-\pi)\,(t_2-t_1)$, la solución es $q_1=0$: se excluye al tipo bajo y al alto se le extrae todo su excedente. Las condiciones de segundo orden se cumplen por la concavidad de $u$.\findemo
\item \emph{Intuición del margen}: reducir un poco $q_1$ cuesta al monopolista, en cada consumidor de tipo bajo, el excedente de la unidad perdida, $t_1u'(q_1)-c$ (ponderado por $\pi$), y le ahorra, en cada consumidor de tipo alto, la renta informativa de esa unidad, $(t_2-t_1)\,u'(q_1)$ (ponderada por $1-\pi$). La condición de primer orden iguala las 2 cosas. Por ejemplo, subir el precio marginal del tipo bajo hace perder algo con los bajos, pero se gana con todos los altos, a los que se puede cobrar más sin que imiten.
\item \emph{Caso habitual en el test}: con $u(q)=[1-(1-q)^2]/2$, $u'(q)=1-q$, la demanda del tipo $i$ es $p=t_i(1-q)$, es decir, $q_i(p)=1-p/t_i$, y su excedente a un precio $p$ es $\mathrm{EC}_i(p)=(t_i-p)^2/2t_i$. Con discriminación perfecta, $p_i=c$ y cuotas $A_i=\mathrm{EC}_i(c)$. Con el menú óptimo, $q_2=1-c/t_2$ y $q_1=1-c\big/\big[t_1-\frac{1-\pi}{\pi}\,(t_2-t_1)\big]$.
\end{itemize}

\subsection{Índice de \autor{Lerner} del monopolio multiproducto}\label{anexo:multiproducto}

\begin{itemize}
\item Con demandas $q_j(\vect{p})$ y costes separables, $B=\sum_jp_j\,q_j(\vect{p})-\sum_jC_j(q_j)$. La condición de primer orden respecto a $p_i$ es
  \begin{equation}
  q_i+\sum_j\big(p_j-\mathrm{CMg}_j\big)\frac{\partial q_j}{\partial p_i}=0 .
  \end{equation}
\item Separando el término $j=i$ y dividiendo por $-p_i\,\partial q_i/\partial p_i>0$:
  \begin{equation}
  \frac{p_i-\mathrm{CMg}_i}{p_i}=\frac{-q_i}{p_i\,\partial q_i/\partial p_i}-\sum_{j\neq i}\big(p_j-\mathrm{CMg}_j\big)\frac{\partial q_j/\partial p_i}{p_i\,\partial q_i/\partial p_i}.
  \end{equation}
\item El primer término es $1/|\varepsilon_{q_i,p_i}|$. En el segundo, $\partial q_j/\partial p_i=\varepsilon_{q_j,p_i}\,q_j/p_i$ y $p_i\,\partial q_i/\partial p_i=\varepsilon_{q_i,p_i}\,q_i=-|\varepsilon_{q_i,p_i}|\,q_i$, así que
  \begin{equation}
  \mathcal{L}_i=\frac{1}{|\varepsilon_{q_i,p_i}|}+\sum_{j\neq i}\frac{\big(p_j-\mathrm{CMg}_j\big)\,q_j}{p_i\,q_i}\cdot\frac{\varepsilon_{q_j,p_i}}{|\varepsilon_{q_i,p_i}|},
  \end{equation}
  que es \eqref{eq:lerner-multi}. El término cruzado lleva la elasticidad de la demanda del bien $j$ respecto al precio del bien $i$.\findemo
\end{itemize}

\subsection{Producción conjunta en proporciones fijas}\label{anexo:proporciones-fijas}

\paragraph{a) Idea}

\begin{itemize}
\item Hay bienes que salen \emph{juntos} de un mismo proceso, en proporciones fijas: la carne y la piel, la lana y la carne del cordero, el coque y el gas. Es la \concepto{oferta conjunta} de \textcite[libro~V, cap.~6]{Marshall1890Principles}, que ya había estudiado \textcite{Mill1848Principles}. No hay un coste propio de cada bien, sino del proceso.
\end{itemize}

\paragraph{b) Modelo}

\subparagraph{i) Supuestos}

\begin{itemize}
\item Cada unidad del proceso, con coste $C(q)$, da una unidad de cada bien.
\item El monopolista puede desechar lo que sobre de cada bien.
\end{itemize}

\subparagraph{ii) Desarrollo}

\begin{itemize}
\item El monopolista elige $q$ y cuánto vende de cada bien, $x_i\leq q$:
  \begin{equation}
  \optim{\max}{q,\,x_1,\,x_2}{p_1(x_1)\,x_1+p_2(x_2)\,x_2-C(q)}{x_1\leq q\\ x_2\leq q}
  \end{equation}
\item Si en el óptimo vende todo lo que produce, el ingreso marginal del proceso es la \emph{suma vertical} de los ingresos marginales: $\mathrm{IMg}_1(q)+\mathrm{IMg}_2(q)=\mathrm{CMg}(q)$.
\item Si a esa cantidad el ingreso marginal de un bien es negativo, le conviene vender de ese bien solo hasta donde su ingreso marginal es nulo (elasticidad unitaria) y \emph{desechar} el resto, y el proceso se ajusta a $\mathrm{IMg}_1(q)=\mathrm{CMg}(q)$. En general, $\sum_i\max\{\mathrm{IMg}_i(q),0\}=\mathrm{CMg}(q)$ (gráfico~\ref{fig:produccion-conjunta}).
\end{itemize}

\begin{figure}[H]
\centering
\caption{Producción conjunta en proporciones fijas: suma vertical de los ingresos marginales}\label{fig:produccion-conjunta}
\grafico[0.6\textwidth]{produccion-conjunta}
\fuente{elaboración propia a partir de \textcite[libro~V, cap.~6]{Marshall1890Principles}.}
\end{figure}

\subparagraph{iii) Implicaciones}

\begin{itemize}
\item No tiene sentido preguntar por el coste de cada bien por separado: el reparto del coste conjunto es arbitrario, y los precios los decide la demanda de cada uno. Un monopolista puede destruir producción perfectamente útil, algo que en competencia no ocurre mientras el precio sea positivo.
\end{itemize}

\subsection{Precios de \autor{Ramsey}-\autor{Boiteux}}\label{anexo:ramsey}

\begin{itemize}
\item Con demandas independientes $q_i(p_i)$, el regulador maximiza $\mathrm{EC}(\vect{p})+B(\vect{p})$ sujeto a $B(\vect{p})\geq 0$. Con el multiplicador $\lambda\geq 0$ de la restricción, maximiza el problema penalizado $\mathrm{EC}(\vect{p})+(1+\lambda)\,B(\vect{p})$, con $\partial\mathrm{EC}/\partial p_i=-q_i$ (la pérdida de excedente de una subida del precio es la cantidad comprada) y $\partial B/\partial p_i=q_i+(p_i-\mathrm{CMg}_i)\,q_i'(p_i)$.
\item Condición de primer orden: $-q_i+(1+\lambda)\big[q_i+(p_i-\mathrm{CMg}_i)\,q_i'\big]=0$, es decir, $(1+\lambda)\,(p_i-\mathrm{CMg}_i)\,q_i'=-\lambda\,q_i$.
\item Dividiendo por $(1+\lambda)\,p_i\,q_i'$:
  \begin{equation}
  \frac{p_i-\mathrm{CMg}_i}{p_i}=\frac{\lambda}{1+\lambda}\cdot\frac{-q_i}{p_i\,q_i'}=\frac{\lambda}{1+\lambda}\cdot\frac{1}{|\varepsilon_{q_i,p_i}|}.
  \end{equation}
  \findemo
\item Si $\lambda\to\infty$, el regulador solo atiende al beneficio y se obtiene la regla del monopolio, $\mathcal{L}_i=1/|\varepsilon_{q_i,p_i}|$; si $\lambda=0$, $p_i=\mathrm{CMg}_i$. El cociente de los márgenes de 2 bienes es inverso al de sus elasticidades: $\mathcal{L}_i/\mathcal{L}_j=|\varepsilon_{q_j,p_j}|/|\varepsilon_{q_i,p_i}|$.
\end{itemize}

\subsection{Sostenibilidad del monopolio natural con reacción del incumbente}\label{anexo:sostenibilidad}

\begin{itemize}
\item En el gráfico~\ref{fig:subaditividad-sostenibilidad}, $\mathrm{CMe}_1$ es el coste medio de una empresa (en U) y $\mathrm{CMe}_{1+2}$ el de 2 empresas que se reparten la producción a partes iguales. Entre el origen y $q^1$ hay subaditividad; entre $q^0$ (mínimo de $\mathrm{CMe}_1$) y $q^1$, sin economías de escala.
\item \emph{Si la demanda corta a $\mathrm{CMe}_1$ antes de $q^0$}, el monopolio natural es deseable y sostenible: para entrar, un competidor tendría que cobrar menos que el incumbente, y en el tramo decreciente de $\mathrm{CMe}_1$ eso es $p<\mathrm{CMe}_1$: tendría pérdidas.
\item \emph{Si la corta entre $q^0$ y $q^1$}, es deseable, pero la sostenibilidad depende de la posición de la demanda frente a $\mathrm{CMe}_{1+2}$:
  \begin{itemize}
  \item \emph{Demanda $D^{*}$, siempre por debajo de $\mathrm{CMe}_{1+2}$}. El incumbente produce $q^{*}$ a $p^{*}$, con beneficio nulo. Si entra otra empresa, con un precio entre $p^0$ y $p^{*}$, el incumbente no puede acomodarla repartiéndose la producción, porque $\mathrm{CMe}_{1+2}$ queda siempre por encima de la demanda y las 2 tendrían pérdidas. Solo le queda luchar bajando el precio (y la cantidad) hasta $p^0$ y $q^0$, con beneficio nulo para él y pérdidas para el entrante. El entrante anticipa esta secuencia y no entra: el monopolio natural es \emph{sostenible} y produce $q^{*}$ a $p^{*}$.
  \item \emph{Demanda $D'$, que en algún tramo supera a $\mathrm{CMe}_{1+2}$}. El incumbente produce $q'$ a $p'$, con beneficio nulo. Si entra otra empresa, con un precio entre $p^0$ y $p'$, el incumbente empieza bajando el precio y la cantidad, pero a partir de cierto punto le compensa acomodar la entrada y repartirse el mercado a un precio por encima de $\mathrm{CMe}_{1+2}$, con beneficios, en lugar de seguir luchando con beneficio nulo ($p=\mathrm{CMe}_1$). El entrante lo anticipa y entra: el monopolio natural \emph{no es sostenible}, y las 2 empresas producen donde el beneficio conjunto es máximo, donde la distancia entre la demanda y $\mathrm{CMe}_{1+2}$ es mayor.\findemo
  \end{itemize}
\item La mayoría de los manuales suponen que el incumbente no reacciona (mantiene su precio y atiende la demanda que deja el entrante) y concluyen que el monopolio natural deja de ser sostenible en cuanto aparecen deseconomías de escala, a partir de $q^0$ \parencite{BaumolPanzarWillig1982Contestable}. Con reacción del incumbente, la frontera está en $\mathrm{CMe}_{1+2}$ \parencite{Segura1986Analisis}. En los 2 casos, la existencia y la sostenibilidad dependen tanto de los costes como de la demanda: la subaditividad es necesaria para la sostenibilidad, y las economías de escala en el tramo relevante bastan para las 2.
\end{itemize}

\subsection{Poder de mercado en el bien y en el trabajo: los 4 casos y el sindicato}\label{anexo:cuatro-casos}

\begin{itemize}
\item Una empresa con poder de monopolio en el bien (demanda inversa $p(q)$) y de monopsonio en el trabajo (oferta inversa $w_L(L)$), con el capital fijo en $\overline{K}$, resuelve
  \begin{equation}
  \optim{\max}{L}{B(L)=p\big(F(L,\overline{K})\big)\,F(L,\overline{K})-w_L(L)\,L-w_K\overline{K}}{L\geq 0}
  \end{equation}
\item La condición de primer orden iguala el ingreso del producto marginal con el gasto marginal:
  \begin{equation}
  \underbrace{\big[p+p'(q)\,q\big]\,\mathrm{PMg}_L}_{\mathrm{IPMg}_L}=\underbrace{w_L(L)+w_L'(L)\,L}_{\mathrm{GMg}_L}\ \Longleftrightarrow\ p\,\mathrm{PMg}_L\left(1-\frac{1}{|\varepsilon_{q,p}|}\right)=w_L(L)\left(1+\frac{1}{\varepsilon_{L,w_L}}\right).
  \end{equation}
\item Con competencia en el bien, $|\varepsilon_{q,p}|\to\infty$ y el primer miembro es $\mathrm{VPMg}_L$; con competencia en el trabajo, $\varepsilon_{L,w_L}\to\infty$ y el segundo es $w_L$. De ahí los 4 casos de la tabla~\ref{tab:cuatro-casos}. El IPMg queda por debajo del VPMg, porque vender más baja el precio de todas las unidades, y el GMg por encima de la oferta, porque contratar más sube el salario de todos. El salario lo da siempre la oferta de trabajo en la cantidad contratada.\findemo
\item \emph{Sindicato monopolista frente a empresas precio-aceptantes}. El sindicato maximiza la renta de sus miembros, $U_S=w\,L-v(L)$, sobre la demanda de trabajo $w^D(L)=\mathrm{VPMg}_L$:
  \begin{equation}
  w^D(L)+w^{D\prime}(L)\,L=v'(L)\ \Longleftrightarrow\ w^D(L)\left(1-\frac{1}{|\varepsilon_{L^D,w}|}\right)=w_L(L).
  \end{equation}
  Su ingreso marginal queda por debajo de la demanda de trabajo, porque ofrecer más trabajo baja el salario de todos. El empleo es menor y el salario mayor que en competencia, y al salario fijado hay más gente dispuesta a trabajar que puestos: hay paro involuntario. Si además la empresa es monopsonista, es el monopolio bilateral (\ref{sec:monopolio-bilateral}).\findemo
\end{itemize}

\subsection{Solución de \autor{Nash} del monopolio bilateral}\label{anexo:nash}

\begin{itemize}
\item Tomando logaritmos en \eqref{eq:nash-bilateral}, se maximiza $\gamma\ln\big[w\,L-v(L)-U_S^0\big]+(1-\gamma)\ln\big[p\,F(L)-w\,L-B^0\big]$. Llamamos $S\equiv p\,F(L)-v(L)-U_S^0-B^0$ al excedente neto que se reparten.
\item Condición de primer orden respecto a $w$: $\dfrac{\gamma\,L}{U_S-U_S^0}=\dfrac{(1-\gamma)\,L}{B-B^0}$; como $(U_S-U_S^0)+(B-B^0)=S$, queda $U_S-U_S^0=\gamma\,S$ y $B-B^0=(1-\gamma)\,S$.
\item Condición de primer orden respecto a $L$: $\dfrac{\gamma\,\big[w-v'(L)\big]}{U_S-U_S^0}+\dfrac{(1-\gamma)\,\big[p\,F'(L)-w\big]}{B-B^0}=0$. Sustituyendo los denominadores, $\big[w-v'(L)\big]+\big[p\,F'(L)-w\big]=0$, es decir,
  \begin{equation}
  p\,F'(L^{*})=v'(L^{*})=w_L(L^{*}):
  \end{equation}
  el empleo maximiza el excedente conjunto, coincide con el competitivo y no depende de $\gamma$.
\item El salario sale de $U_S-U_S^0=\gamma\,S$: $w^{*}L^{*}=v(L^{*})+U_S^0+\gamma\,\big[p\,F(L^{*})-v(L^{*})-U_S^0-B^0\big]$.\findemo
\end{itemize}

\subsection{Efecto \autor{Averch}-\autor{Johnson}}\label{anexo:aj}

\begin{itemize}
\item Sea $I(z_K,z_L)=p\big(F(z_K,z_L)\big)\,F(z_K,z_L)$ el ingreso, con $I_i=\mathrm{IMg}\cdot\mathrm{PMg}_i$. Con el multiplicador $\mu$ de la restricción de rentabilidad, el problema penalizado es $I-w_Kz_K-w_Lz_L-\mu\,\big[I-w_Lz_L-s\,z_K\big]$; con la restricción activa, $0<\mu<1$.
\item Condición de primer orden respecto a $z_L$: $(1-\mu)\,I_L-(1-\mu)\,w_L=0\Rightarrow I_L=w_L$.
\item Condición de primer orden respecto a $z_K$: $(1-\mu)\,I_K-w_K+\mu\,s=0\Rightarrow I_K=w_K-\dfrac{\mu}{1-\mu}\,(s-w_K)$.
\item Dividiendo, $\dfrac{\mathrm{PMg}_L}{\mathrm{PMg}_K}=\dfrac{I_L}{I_K}=\dfrac{w_L}{w_K-\frac{\mu}{1-\mu}(s-w_K)}$. Con $s>w_K$ y $0<\mu<1$, el denominador es menor que $w_K$: la empresa usa capital hasta donde su productividad marginal es menor que la que exige el mínimo coste. Sobrecapitaliza.\findemo
\end{itemize}

% ==========================================================================
\section{Anexo: aspectos para el test y preguntas del tribunal}\label{anexo:test}

\begin{anotaciones}
Preguntas tipo test del primer ejercicio: \href{https://www.victorgutierrezmarcos.es/oposicion/temario/primer-ejercicio/test/simulador.html}{simulador de test}.
\end{anotaciones}

\begin{itemize}
\item \emph{Demanda lineal} $p=a-bq$ y CMg constante $c$: $\mathrm{IMg}=a-2bq$ (misma ordenada en el origen y doble pendiente); $q^M=(a-c)/2b$, $p^M=(a+c)/2$, $B=(a-c)^2/4b$ sin costes fijos, y $\mathrm{PIE}=(a-c)^2/8b=B/2$. Con un impuesto unitario $t$, $p^M$ sube $t/2$.
\item \emph{Demanda isoelástica} $q=A\,p^{-\eta}$, con $|\varepsilon_{q,p}|=\eta>1$, y CMg constante $c$: $p^M=\frac{\eta}{\eta-1}\,c$. Ejemplo: con $q=p^{-2}$ y $C(q)=5q+10$, $p^M(1-1/2)=5$, así que $p^M=10$ y $q^M=1/100$ (en competencia, $p=5$ y $q=1/25$). Con ese coste fijo el beneficio es negativo, $B=0{,}1-0{,}05-10<0$: el monopolista solo produciría a corto plazo, porque $p^M>\mathrm{CVMe}=5$. Se obtiene igual con $B(q)=\sqrt{q}-5q-10$, con $\mathrm{IMg}=\mathrm{CMg}$ o con la regla del margen. Con demanda isoelástica, un impuesto unitario sube el precio en más que el impuesto.
\item \emph{Tramo elástico}: el monopolista nunca produce donde $|\varepsilon_{q,p}|\leq 1$ si $\mathrm{CMg}>0$.
\item \emph{Tarifa en dos partes con un solo tipo} ($n$ consumidores iguales): $p=c$ y $A=\mathrm{EC}(c)$; el monopolista se queda con todo el excedente, como en el primer grado.
\item \emph{Tarifa en dos partes común con 2 tipos} ($n_1$ bajos, $n_2$ altos), sin excluir a nadie: $A=\mathrm{EC}_1(p)$ y $B=(n_1+n_2)\,A+(p-c)\,[n_1q_1(p)+n_2q_2(p)]$, que se maximiza en $p$ (con 2 consumidores, $B=2A+(p-c)(q_1+q_2)$); el resultado es $p>c$. Hay que compararlo con la alternativa de excluir a los bajos: $p=c$ y $A=\mathrm{EC}_2(c)$, con beneficio $n_2\,\mathrm{EC}_2(c)$.
\item \emph{Segundo grado con 2 tipos}: no hay distorsión arriba; el tipo bajo compra menos de lo eficiente y no tiene excedente; el alto se queda con una renta informativa.
\item \emph{Tercer grado}: el precio es mayor en el mercado de demanda menos elástica; con demandas lineales y todos los mercados servidos, la cantidad total no cambia respecto al precio uniforme y el excedente total baja.
\item \emph{Primer grado}: eficiente (la cantidad es la competitiva), excedente del consumidor nulo.
\item \emph{Monopolio natural}: la subaditividad es la definición; las economías de escala son suficientes pero no necesarias; las economías de alcance, necesarias pero no suficientes en el multiproducto; las pruebas de subsidios cruzados, necesarias para la sostenibilidad.
\item \emph{Ramsey-Boiteux}: el margen relativo es proporcional a la inversa de la elasticidad; con $\lambda=0$, $p=\mathrm{CMg}$, y el margen crece con $\lambda$.
\item \emph{Monopsonio con oferta lineal} $w_L=\alpha+\beta L$: $\mathrm{GMg}_L=\alpha+2\beta L$ (misma ordenada y doble pendiente, como el IMg de una demanda lineal). Con monopsonio no hay paro involuntario; con un salario mínimo entre $w^M$ y $w^c$ el empleo sube.
\item \emph{Doble marginalización}: con demanda lineal, la cadena separada vende la mitad que la integrada.
\item \emph{Fechas}: \autor{Cournot} (1838), \autor{Dupuit} (1844, 1849), \autor{Marshall} (1890), \autor{Pigou} (1920), \autor{Bowley} (1928), \autor{Robinson} (1933), \autor{Lerner} (1934), \autor{Coase} (1946 y 1972), \autor{Harberger} (1954), \autor{Boiteux} (1956), \autor{Averch} y \autor{Johnson} (1962), \autor{Tullock} (1967), \autor{Demsetz} (1968), \autor{Stigler} (1971), \autor{Posner} (1975), \autor{Faulhaber} (1975), \autor{Baumol} (1977), \autor{Mussa} y \autor{Rosen} (1978), \autor{Littlechild} (1983), \autor{Shleifer} (1985), \autor{Laffont} y \autor{Tirole} (1986, 1993).
\end{itemize}

\paragraph{Preguntas probables del tribunal}

\begin{itemize}
\item ¿Por qué el monopolista no tiene curva de oferta? ¿Por qué nunca produce en el tramo inelástico?
\item ¿Es la pérdida de bienestar del monopolio solo el triángulo de \autor{Harberger}? ¿Cuánto vale con demanda lineal?
\item ¿Es eficiente la discriminación de primer grado? ¿Cuándo aumenta el excedente total la de tercer grado?
\item ¿Por qué en la discriminación de segundo grado se distorsiona la cantidad del tipo bajo y no la del alto?
\item ¿Son las economías de escala necesarias para que haya monopolio natural? ¿Y las de alcance en el multiproducto?
\item ¿Qué dice la conjetura de \autor{Coase}?
\item En la regla de \autor{Ramsey}-\autor{Boiteux}, ¿qué ocurre con los márgenes si aumenta $\lambda$?
\item ¿Hay paro involuntario con monopsonio? ¿Por qué un salario mínimo puede aumentar el empleo?
\item En el monopolio bilateral, ¿depende la cantidad del poder de negociación?
\item ¿Qué es el efecto \autor{Averch}-\autor{Johnson}? ¿Qué ventajas tiene el \emph{price cap} frente a la regulación de la tasa de rentabilidad?
\end{itemize}

\nocite{ICEXCECO3A17,Nobel2014TiroleLecture,ViscusiHarringtonVernon2005Economics}
\bibliografia

\end{document}
